% ============================================================================
% Section: Main results -- disorder-free localization
% ============================================================================
% Section-local macros (unique names, \providecommand only).
\providecommand{\cdr}{\tilde c}
\providecommand{\bent}{\mathsf{h}}           % binary entropy function (not the field h)
\providecommand{\Hsh}{\mathrm{H}}            % Shannon entropy
\providecommand{\Sfm}{S_{f}}                 % bare-matter entropy
\providecommand{\Ibar}{\bar I}
\providecommand{\rhoqf}[1]{\rho[#1]}         % quasi-free density matrix with symbol #1
\providecommand{\Alg}{\mathfrak{A}}
\providecommand{\BAlg}{\mathfrak{B}}
\providecommand{\GAlg}{\mathfrak{G}}
\providecommand{\SigF}{\Sigma_1(F)}
\providecommand{\SigpF}{\Sigma_1'(F)}
\providecommand{\ThF}{\Theta(F)}
\providecommand{\FZ}{F_{\Z}}
\providecommand{\Zone}{\one_N}

\section{Main results: disorder-free localization}\label{sec:main}

By the duality (\cref{thm:duality}) and the exact sector averaging (\cref{cor:sector}), every
expectation value of gauge-invariant observables (at any number of times) in the initial state
$\Psi_0=\ket{\uparrow\cdots\uparrow}_z\otimes\psi$ is an average over $2^N$ independent free-fermion
problems, one for each charge configuration $q$, with the i.i.d.\ Bernoulli one-body operators
$\hq=-J\Delta_N+2h\,\diag(q)$.  In this section we turn annealed one-body bounds into
statements about the gauge theory.

\emph{Status of the results.}  In finite volume all \emph{localization bounds} are stated under
Hypothesis~\ref{hyp:EC}, with explicit constants in terms of the moments $M_p(F)$, $M'_p(F)$ of
the function $F$ in \eqref{eq:hypEC}.  The structural results need no hypothesis: the lemmas of
\S\ref{sec:main:prelim} other than \cref{lem:annealed}, \cref{prop:LR}(a),(c),(d),
\cref{thm:imbalance}(i)--(iii), \cref{thm:arealaw}(a), \cref{prop:entropies}, \cref{lem:intentropy}
and the first bound of \cref{lem:duhamel}.
For every $\lambda=2h/J\neq0$, Theorem~\ref{thm:EC} proves Hypothesis~\ref{hyp:EC} with an exponential
function, $F(d)=\min\{1,Ce^{-\eta d}\}$.  So the finite-volume results hold unconditionally, but $C$ and $\eta$
are not explicit (\cref{rem:rate}).  The infinite-volume
statements (\cref{thm:infvol}) use only the Ge--Zhao bound on $\Z$ (\cref{prop:Zline}) and are
independent of Theorem~\ref{thm:EC}.  The constants are far from optimal; we state them because they
are explicit.  Two statements are conjectures (Conjectures~\ref{conj:logent} and~\ref{conj:volmix}).
Several arguments adapt methods developed for disordered XY chains
\cite{HSS2012,AbdulRahmanStolz2015,ANSS2016}, with the sector average in place of the average over a random
field; we say so where this applies (\cref{rem:main-transport,rem:main-LR,rem:main-arealaw}).

\emph{Overview.}
\Cref{thm:transport}: sector-averaged localization of the dressed-fermion propagator, bounded
moments and bounded transport across every cut.
\Cref{prop:LR}: a sector-averaged zero-velocity Lieb--Robinson bound, and the proof that no
operator-norm version exists.
\Cref{thm:imbalance}: a strictly positive infinite-time density-wave imbalance for every $h\neq0$,
uniformly in $N$.
\Cref{thm:arealaw,prop:entropies}: uniform-in-time area laws, and the precise relations between the
entropies $\SG,\Sfull,\Smix,\Sfm,\SN$.
\Cref{thm:numberentropy}: a uniform bound on the number entropy.
\Cref{thm:infvol}: thermodynamic limit.
\Cref{lem:duhamel} and Conjectures~\ref{conj:logent} and~\ref{conj:volmix}: what lies beyond.

\subsection{Preliminaries}\label{sec:main:prelim}

\paragraph{Standing notation.}
Throughout, $J>0$, $h\in\R\setminus\{0\}$ unless said otherwise, $N\ge2$, and $\E=\E_q$ is the
expectation over $q$ uniform on $\{\pm1\}^N$ (in finite volume a finite average, so no measurability
questions arise).  We write
\[
  U_q(t):=e^{-it\hq},\qquad Q_q:=Q^{[1,N]}_q,
\]
and $R_q$ for the time-averaged transition probability \eqref{eq:Rdef} of $\hq$.  Since $\hq$ is a
real symmetric matrix, $U_q(t)^{\top}=U_q(t)$, so $|U_q(t)_{xy}|=|U_q(t)_{yx}|$; and
$\sum_x|U_q(t)_{xy}|^2=1$.  The dressed fermions $\cdr_x$ are those of \eqref{eq:dressed}, and
$O(t):=e^{itH}Oe^{-itH}$ for operators $O$ on $\cH_N$.  For a Fock state $\psi=\ket{n}$ we write
$S:=\{x:n_x=1\}$ for its occupied set and
\begin{equation}\label{eq:main-K}
  K_q(t):=U_q(t)P_SU_q(t)^*,
\end{equation}
where $P_S$ is the orthogonal projection onto $\ell^2(S)$.  $K_q(t)$ is the projection onto
$U_q(t)\ell^2(S)$, and by \cref{cor:sector}(c),
$\langle\psi_q(t),c_x^\dagger c_y\psi_q(t)\rangle=\sum_{s\in S}\overline{U_q(t)_{xs}}U_q(t)_{ys}=K_q(t)_{yx}$.
The region $A$ consists of the sites $1,\dots,\ell$ and the links $b_1,\dots,b_\ell$ ($1\le\ell<N$), and
$\rho^q_A(t)$ denotes the reduced density matrix of $\psi_q(t)$ on the modes $1,\dots,\ell$, with the
identification of \S\ref{sec:model:cut}.  $\Hsh(p)=-\sum_kp(k)\ln p(k)$ is the Shannon entropy and
$\bent(\nu):=-\nu\ln\nu-(1-\nu)\ln(1-\nu)$ the binary entropy function (not to be confused with the
field $h$).

\paragraph{Constants.}  For $F$ as in Hypothesis~\ref{hyp:EC} we use
\begin{equation}\label{eq:main-constants}
\begin{gathered}
  \SigF:=\sum_{d\ge1}d\,F(d)\le M_1(F),\qquad
  \SigpF:=\sum_{d\ge1}d\,F(d)^{1/2}\le M'_1(F),\\
  \ThF:=2\sum_{k\ge1}k(2k+1)F(k)\le4M_2(F).
\end{gathered}
\end{equation}
The last inequality holds because $2k(2k+1)\le4(k+1)^2$.

\begin{lemma}[Sector-diagonal operators]\label{lem:sectornorm}
Let $X$ be a gauge-invariant operator on $\cH_N$.
\begin{enumerate}[label=(\alph*)]
\item $X=\sum_qP_qXP_q$ and $\norm{X}=\max_q\norm{P_qXP_q}$.
\item Let $\rho$ be a density matrix on $\cH_N$ and $w_q:=\Tr(\rho P_q)$ its sector weights.  Then
$|\Tr(\rho X)|\le\sum_qw_q\norm{P_qXP_q}$.  One has $w_q=2^{-N}$, and hence
$|\Tr(\rho X)|\le\E_q\norm{P_qXP_q}$, for $\rho=4^{-N}\one$ and for $\rho=\ket{\Psi_0}\bra{\Psi_0}$ with
$\Psi_0=\ket{\uparrow\cdots\uparrow}_z\otimes\psi$, where $\psi\in\cF_N$ is any unit vector.
\item If $X=\sum_qx_qP_q$ with $x_q\in\C$, then
$\langle\Psi_0,|X|\Psi_0\rangle=4^{-N}\Tr|X|=\E_q|x_q|$ for every $\Psi_0$ as in (b).
\end{enumerate}
\end{lemma}

\begin{proof}
(a) $X$ commutes with the orthogonal projections $P_q$, which sum to $\one$; the blocks $P_qXP_q$ act on
mutually orthogonal subspaces.  (b) By (a), $\Tr(\rho X)=\sum_q\Tr(P_q\rho P_q\,P_qXP_q)$, and
$|\Tr(\sigma Y)|\le\Tr(\sigma)\norm{Y}$ for $\sigma\ge0$.  The weights are $2^{-N}$ by
\cref{cor:sector}(a), resp.\ because $\Tr P_q=2^N$ (\cref{lem:gauss}(iii)).  (c) $|X|=\sum_q|x_q|P_q$; use
$\norm{P_q\Psi_0}^2=2^{-N}=4^{-N}\Tr P_q$.
\end{proof}

\paragraph{Quasi-free density matrices.}
\Cref{lem:quasifree,lem:slater} below are standard facts about quasi-free states.  We include short
proofs because we use them for arbitrary, also non-contiguous, sets of modes, with an explicit
identification of the tensor factors.
Let $A$ be a finite set of fermionic modes, $\cF_A$ the Fock space over $\ell^2(A)$, $c_x$ ($x\in A$) the
annihilation operators, $c(v):=\sum_x\overline{v(x)}c_x$, and $\Gamma(C)$ the second quantization of a
matrix $C$ on $\ell^2(A)$ ($\Gamma(C)\ket0=\ket0$,
$\Gamma(C)c^\dagger(u_1)\cdots c^\dagger(u_k)\ket0=c^\dagger(Cu_1)\cdots c^\dagger(Cu_k)\ket0$).  For
$0\le C\le\one$ with eigenvalues $\nu_1,\dots,\nu_{|A|}$ and an orthonormal eigenbasis $(v_i)$, put
$d_i:=c(v_i)$ and
\begin{equation}\label{eq:main-rhoC}
  \rhoqf{C}:=\prod_{i=1}^{|A|}\bigl(\nu_i\,d_i^\dagger d_i+(1-\nu_i)\,d_id_i^\dagger\bigr)
\end{equation}
(the factors commute).

\begin{lemma}[Quasi-free density matrices]\label{lem:quasifree}
Let $0\le C\le\one$ on $\ell^2(A)$.
\begin{enumerate}[label=(\alph*)]
\item With $N_A:=\sum_{x\in A}c_x^\dagger c_x$ and the convention that the $j=0$ term is $\Gamma(C)$,
\begin{equation}\label{eq:main-rhoformula}
  \rhoqf{C}=\sum_{j=0}^{|A|}\frac{(-1)^j}{j!}\sum_{x_1,\dots,x_j\in A}c_{x_1}\cdots c_{x_j}\,\Gamma(C)\,
  c^\dagger_{x_j}\cdots c^\dagger_{x_1}.
\end{equation}
In particular $\rhoqf{C}$ does not depend on the choice of eigenbasis, its matrix elements in the
occupation basis are polynomials in the entries of $C$, and $C\mapsto\rhoqf{C}$ is uniformly continuous
on $\{0\le C\le\one\}$.
\item $\rhoqf{C}$ is a density matrix with $S(\rhoqf{C})=\sum_i\bent(\nu_i)$.  Under $\rhoqf{C}$, $N_A$ has
the law of $\sum_i\xi_i$ with independent $\xi_i\sim\mathrm{Bernoulli}(\nu_i)$; its generating function is
$\Tr(\rhoqf{C}z^{N_A})=\det(\one+(z-1)C)$, and its variance is $\Tr(C-C^2)$.
\end{enumerate}
\end{lemma}

\begin{proof}
(a) For every operator $Y$ on $\cF_A$,
$\sum_id_iYd_i^\dagger=\sum_{x,y}\bigl(\sum_i\overline{v_i(x)}v_i(y)\bigr)c_xYc_y^\dagger=\sum_xc_xYc_x^\dagger$,
by completeness of $(v_i)$.  Iterating, the inner sum in \eqref{eq:main-rhoformula} equals
$\Xi_j:=\sum_{i_1,\dots,i_j}d_{i_1}\cdots d_{i_j}\Gamma(C)d^\dagger_{i_j}\cdots d^\dagger_{i_1}$.  Let
$\ket{T}:=\prod_{i\in T}d_i^\dagger\ket0$ (increasing order) for $T\subset\{1,\dots,|A|\}$.  The vector
$\phi:=d^\dagger_{i_j}\cdots d^\dagger_{i_1}\ket T$ vanishes unless $i_1,\dots,i_j$ are distinct and not in
$T$; then $\phi=\pm\ket{T\cup R}$ with $R=\{i_1,\dots,i_j\}$, $\Gamma(C)\phi=\prod_{i\in T\cup R}\nu_i\,\phi$
because $Cv_i=\nu_iv_i$, and $d_{i_1}\cdots d_{i_j}\phi=\ket T$.  Each set $R$ arises from $j!$ ordered
tuples, so
\[
  \sum_j\frac{(-1)^j}{j!}\Xi_j\ket T=\sum_{R\subset T^c}(-1)^{|R|}\prod_{i\in T\cup R}\nu_i\,\ket T
  =\prod_{i\in T}\nu_i\prod_{i\notin T}(1-\nu_i)\,\ket T=\rhoqf{C}\ket T .
\]
The right side of \eqref{eq:main-rhoformula} involves only $C$, and the matrix elements of $\Gamma(C)$ in
the occupation basis are minors of $C$.  A continuous function on a compact set is uniformly continuous.

(b) In the basis $(\ket T)$, $\rhoqf{C}$ is diagonal with eigenvalues
$w_T=\prod_{i\in T}\nu_i\prod_{i\notin T}(1-\nu_i)\ge0$, the law of independent Bernoulli$(\nu_i)$
variables, and $N_A=\sum_id_i^\dagger d_i$ (completeness again) acts on $\ket T$ as $|T|$.  The entropy of a
product law is the sum of the entropies of its factors.  Finally
$\sum_Tw_Tz^{|T|}=\prod_i(1-\nu_i+z\nu_i)=\det(\one+(z-1)C)$, and
$\Var\sum_i\xi_i=\sum_i\nu_i(1-\nu_i)$.
\end{proof}

\begin{lemma}[Reduced states of Slater determinants]\label{lem:slater}
Let $\Lambda$ be a finite set of modes, $W\subset\ell^2(\Lambda)$ a subspace of dimension $M$ with orthogonal
projection $K$, and $\psi_W:=c^\dagger(w_1)\cdots c^\dagger(w_M)\ket0$ for an orthonormal basis $(w_j)$ of $W$
(unique up to a phase).  Let $A\subset\Lambda$, $B:=\Lambda\setminus A$, and identify
$\cF_\Lambda\cong\cF_A\otimes\cF_B$ by $c(a)\mapsto c(a)\otimes\one$, $c(b)\mapsto(-1)^{N_A}\otimes c(b)$ for
$a\in\ell^2(A)$, $b\in\ell^2(B)$, $\ket0\mapsto\ket0\otimes\ket0$.  (For $A=\{1,\dots,\ell\}$ this is the
identification of \S\ref{sec:model:cut}.)  Under it, the CAR algebra of the modes in $A$ is
$\mathcal B(\cF_A)\otimes\one$, and we let $\rho_A:=\Tr_{\cF_B}\ket{\psi_W}\bra{\psi_W}$.  Let
$C_A:=P_AKP_A$, regarded as an operator on $\ell^2(A)$, with eigenvalues $\nu_1,\dots,\nu_{|A|}$, and
$T:=P_BKP_A$.  Then:
\begin{enumerate}[label=(\alph*)]
\item $\rho_A=\rhoqf{C_A}$.  Hence $S(\rho_A)=\sum_i\bent(\nu_i)$, and the law of $N_A$ under $\psi_W$ is the
Poisson-binomial law of \cref{lem:quasifree}(b).
\item $C_A-C_A^2=T^*T$.  Hence
$\Var_{\psi_W}(N_A)=\norm{T}_2^2=\sum_{x\in A,\,y\in B}|K_{yx}|^2$ and
$\sum_i\sqrt{\nu_i(1-\nu_i)}=\norm{T}_1\le\sum_{x\in A,\,y\in B}|K_{yx}|$.
\end{enumerate}
\end{lemma}

\begin{proof}
(a) Changing the basis of $W$ multiplies $\psi_W$ by a determinant of modulus one, so we may take $w_j$ to be
an orthonormal eigenbasis of $W$ for the operator $(P_A|_W)^*(P_A|_W)$, with eigenvalues
$\mu_j\in[0,1]$.  Then $\langle P_Aw_j,P_Aw_k\rangle=\mu_j\delta_{jk}$ and
$\langle P_Bw_j,P_Bw_k\rangle=(1-\mu_j)\delta_{jk}$.  Put $a_j:=\mu_j^{-1/2}P_Aw_j$ if $\mu_j>0$ and
$b_j:=(1-\mu_j)^{-1/2}P_Bw_j$ if $\mu_j<1$; these are orthonormal families in $\ell^2(A)$, resp.\ $\ell^2(B)$,
and $w_j=\sqrt{\mu_j}\,a_j+\sqrt{1-\mu_j}\,b_j$ (a term with vanishing coefficient is absent).  Expanding
the product $\prod_jc^\dagger(w_j)$ and moving all $c^\dagger(a_j)$ to the left,
\[
  \psi_W=\sum_{T\subset\{1,\dots,M\}}\varepsilon_T\prod_{j\in T}\sqrt{\mu_j}\prod_{j\notin T}\sqrt{1-\mu_j}\;
  c^\dagger(a_T)\,c^\dagger(b_{T^c})\ket0,
\]
with signs $\varepsilon_T=\pm1$ and $c^\dagger(a_T):=\prod_{j\in T}c^\dagger(a_j)$ in increasing order.  Under the
identification, $c^\dagger(a_T)c^\dagger(b_{T^c})\ket0\mapsto\alpha_T\otimes\beta_{T^c}$ with
$\alpha_T:=c^\dagger(a_T)\ket0$, $\beta_{T^c}:=c^\dagger(b_{T^c})\ket0$, because the factors $(-1)^{N_A}$ act on the
vacuum of $\cF_A$.  The $\alpha_T$ with nonzero coefficient are orthonormal, and so are the $\beta_{T^c}$.  Hence
this is a Schmidt decomposition, and
$\rho_A=\sum_T\prod_{j\in T}\mu_j\prod_{j\notin T}(1-\mu_j)\,\ket{\alpha_T}\bra{\alpha_T}$.  On the other hand
$C_A=\sum_j\ket{P_Aw_j}\bra{P_Aw_j}=\sum_{j:\mu_j>0}\mu_j\ket{a_j}\bra{a_j}$, so $C_A$ has the eigenvectors
$a_j$ ($\mu_j>0$) with eigenvalues $\mu_j$, and eigenvalue $0$ on the orthogonal complement in $\ell^2(A)$.
Computing \eqref{eq:main-rhoC} in this eigenbasis gives the same operator.  The rest is
\cref{lem:quasifree}(b); $N_A\otimes\one$ belongs to the algebra of $A$.

(b) Since $K=K^*=K^2$, $C_A-C_A^2=P_AK(\one-P_A)KP_A=T^*T$.  So
$\sum_i\nu_i(1-\nu_i)=\Tr T^*T=\norm{T}_2^2$ and $\sum_i\sqrt{\nu_i(1-\nu_i)}=\Tr|T|=\norm T_1$.  Finally
$T=\sum_{x\in A,y\in B}K_{yx}\ket{\delta_y}\bra{\delta_x}$ is a sum of rank-one operators of trace norm
$|K_{yx}|$.
\end{proof}

\begin{lemma}[Binary entropy]\label{lem:binent}
For all $\nu\in[0,1]$: $\bent(\nu)\le2\ln2\,\sqrt{\nu(1-\nu)}$, with equality iff $\nu\in\{0,\tfrac12,1\}$.
\end{lemma}

\begin{proof}
Both sides vanish at $\nu\in\{0,1\}$ and are symmetric under $\nu\mapsto1-\nu$.  Write $\nu=(1-s)/2$ with
$s\in[0,1)$.  The function $f(s):=(1+s)\ln(1+s)+(1-s)\ln(1-s)$ satisfies $f(0)=f'(0)=0$ and
$f''(s)=2/(1-s^2)=2\sum_{j\ge0}s^{2j}$, hence $f(s)=\sum_{k\ge1}s^{2k}/(k(2k-1))$ and
\[
  \bent(\nu)=\ln2-\tfrac12f(s)=\ln2-\sum_{k\ge1}e_ks^{2k},\qquad e_k:=\frac1{2k(2k-1)} .
\]
Moreover $2\ln2\sqrt{\nu(1-\nu)}=\ln2\sqrt{1-s^2}=\ln2\,\bigl(1-\sum_{k\ge1}c_ks^{2k}\bigr)$ with
$c_k:=\binom{2k}{k}\frac1{(2k-1)4^k}>0$.  So the claim is $\sum_{k\ge1}d_ks^{2k}\ge0$ with
$d_k:=e_k-c_k\ln2$.  We have $d_1=\frac12(1-\ln2)>0$.  The ratio
$e_k/c_k=4^k/\bigl(2k\binom{2k}{k}\bigr)$ satisfies $(e_{k+1}/c_{k+1})/(e_k/c_k)=2k/(2k+1)<1$, and
$e_2/c_2=\frac23<\ln2$; hence $d_k<0$ for all $k\ge2$.  Letting $s\uparrow1$ (monotone convergence; all
$e_k,c_k>0$) gives $\sum_ke_k=\frac12f(1)=\ln2$ and $\sum_kc_k=1$, so $\sum_{k\ge2}d_k=-d_1$.  Since
$s^{2k}\le s^4$ for $k\ge2$,
\[
  \sum_{k\ge1}d_ks^{2k}\ge d_1s^2+s^4\sum_{k\ge2}d_k=d_1(s^2-s^4)\ge0,
\]
with equality for $s\in[0,1)$ only if $s=0$, i.e.\ $\nu=\frac12$.
\end{proof}

\begin{lemma}[Annealed one-body bounds]\label{lem:annealed}
Assume Hypothesis~\ref{hyp:EC}.  Let $\psi$ be a Fock state with occupied set $S$ (arbitrary), and
$x\neq y$ in $[1,N]$, $d:=|x-y|$.
\begin{enumerate}[label=(\roman*)]
\item $\displaystyle\E\sup_{t\in\R}|K_q(t)_{xy}|\le\E\sum_{z=1}^NQ_q(x,z)Q_q(y,z)\le2\sum_{k\ge\lceil d/2\rceil}F(k)$.
\item For $1\le m<N$:
\begin{gather*}
  \E\sup_t\sum_{x\le m<y}|K_q(t)_{xy}|\le\ThF,\qquad
  \E\sum_{x\le m<z}Q_q(x,z)\le\SigF,\\
  \E\Bigl(\sum_{x\le m<z}Q_q(x,z)\Bigr)^2\le\SigpF^2 .
\end{gather*}
\item For an interval $B\subset[1,N]$ the same bounds hold for the sums over pairs $x\in B$,
$y\in[1,N]\setminus B$, with the constants $2\ThF$, $2\SigF$ and $4\SigpF^2$.
\end{enumerate}
\end{lemma}

\begin{proof}
(i) $|K_q(t)_{xy}|\le\sum_{s\in S}|U_q(t)_{xs}||U_q(t)_{ys}|\le\sum_z Q_q(x,z)Q_q(y,z)$ for all $t$, by
\cref{lem:consequences}(b).  By \eqref{eq:cons2},
$\E[Q_q(x,z)Q_q(y,z)]\le F(\max\{|x-z|,|y-z|\})$.  Let $x<y$.  For $z\in[x,y]$ the value
$k=\max\{z-x,y-z\}$ lies in $[\lceil d/2\rceil,d]$, and each such $k$ arises for at most two $z$; for $z<x$
or $z>y$ it equals $d+|z-x|$, resp.\ $d+|z-y|$, so each $k>d$ arises at most twice.  Hence the sum is at
most $2\sum_{k\ge\lceil d/2\rceil}F(k)$.

(ii) The number of pairs $x\le m<y$ with $y-x=d$ is at most $d$.  By (i),
\[
  \E\sup_t\sum_{x\le m<y}|K_q(t)_{xy}|\le\sum_{d\ge1}d\cdot2\sum_{k\ge\lceil d/2\rceil}F(k)
  =2\sum_{k\ge1}F(k)\sum_{d=1}^{2k}d=2\sum_{k\ge1}k(2k+1)F(k).
\]
The second bound is \eqref{eq:hypEC} with the same counting.  For the third, by \eqref{eq:cons3},
$\E(\sum_iQ_i)^2=\sum_{i,j}\E Q_iQ_j\le(\sum_iF(d_i)^{1/2})^2\le\SigpF^2$, where $i$ runs over the pairs
and $d_i$ is the distance of pair $i$.

(iii) The number of pairs $x\in B$, $y\notin B$ at distance $d$ is at most $2d$.
\end{proof}

\begin{lemma}[Finite propagation speed]\label{lem:propagation}
Let $\Lambda\subset\Lambda'\subset\Z$ be intervals ($\Lambda'=\Z$ allowed), $v:\Lambda'\to\R$ with $|v|\le V_0$,
$k_{\Lambda}:=-J\Delta_\Lambda+\diag(v|_\Lambda)$ on $\ell^2(\Lambda)$, likewise $k_{\Lambda'}$, and $c:=|t|(2J+V_0)$.  For
$x,y\in\Lambda$:
\begin{enumerate}[label=(\alph*)]
\item $|(e^{-itk_\Lambda})_{xy}|\le e^{c}\,c^{|x-y|}/|x-y|!$;
\item $|(e^{-itk_\Lambda})_{xy}-(e^{-itk_{\Lambda'}})_{xy}|\le2e^{c}\,c^{n_0}/n_0!$, where
$n_0:=\dist(x,\Lambda'\setminus\Lambda)+\dist(y,\Lambda'\setminus\Lambda)$ (the left side is $0$ if $\Lambda=\Lambda'$).
\end{enumerate}
\end{lemma}

\begin{proof}
$\norm{k_\Lambda},\norm{k_{\Lambda'}}\le2J+V_0$, and $e^{-itk}=\sum_n(-it)^nk^n/n!$ converges in norm.  The
entry $(k^n)_{xy}$ is a sum over paths $x=z_0,z_1,\dots,z_n=y$ with $|z_{i+1}-z_i|\le1$ of products of matrix
entries.  (a) $(k_\Lambda^n)_{xy}=0$ for $n<|x-y|$ and $|(k^n_\Lambda)_{xy}|\le\norm{k_\Lambda}^n$, so the entry is at
most $\sum_{n\ge d}c^n/n!\le e^cc^d/d!$ with $d=|x-y|$.  (b) A path from $x$ to $y$ that visits
$\Lambda'\setminus\Lambda$ has length at least $n_0$, and on paths inside $\Lambda$ the entries of $k_\Lambda$ and
$k_{\Lambda'}$ agree.  So $(k_\Lambda^n)_{xy}=(k^n_{\Lambda'})_{xy}$ for $n<n_0$, and the difference is at most
$2\sum_{n\ge n_0}c^n/n!\le2e^cc^{n_0}/n_0!$.
\end{proof}

\subsection{Localization of dressed fermions and bounded transport}\label{sec:main:transport}

\begin{theorem}[Dressed-fermion localization and bounded transport]\label{thm:transport}
Assume Hypothesis~\ref{hyp:EC}, and let $\Psi_0=\ket{\uparrow\cdots\uparrow}_z\otimes\psi$,
$\Psi(t)=e^{-itH}\Psi_0$.  For all $N\ge2$:
\begin{enumerate}[label=(\alph*)]
\item \emph{(Propagator.)}  For all $x,y\in[1,N]$ and $t\in\R$,
\begin{equation}\label{eq:main-anticomm}
  \{\cdr_x(t),\cdr_y^\dagger\}=\sum_qU_q(t)_{xy}\,P_q .
\end{equation}
Thus the dressed-fermion propagator is the scalar $U_q(t)_{xy}$ in each sector, and its sector-averaged
modulus is localized uniformly in time:
\begin{equation}\label{eq:main-prop}
  \E_q\sup_{t\in\R}\norm{P_q\{\cdr_x(t),\cdr_y^\dagger\}P_q}=\E_q\sup_{t\in\R}|U_q(t)_{xy}|\le F(|x-y|).
\end{equation}
In particular, $\sup_t\langle\Psi_0,|\{\cdr_x(t),\cdr_y^\dagger\}|\,\Psi_0\rangle\le F(|x-y|)$ for every
matter state $\psi$.
\item \emph{(One particle: moments.)}  If $\psi=f_y^\dagger\ket0$, then for $0\le p\le2$,
\[
  \sup_{t\in\R}\sum_{x=1}^N(|x-y|+1)^p\langle\Psi(t),n_x\Psi(t)\rangle
  \le\E_q\sup_{t\in\R}\sum_{x=1}^N(|x-y|+1)^p|U_q(t)_{xy}|^2\le2M_p(F).
\]
\item \emph{(Transport across any cut.)}  Let $\psi=\ket n$ be a Fock state with occupied set $S$, let
$1\le m<N$ and $s_m:=|S\cap[1,m]|$.  Then for every $q$ and $t$,
$|\langle\psi_q(t),N_{\le m}\psi_q(t)\rangle-s_m|\le\sum_{x\le m<z}|U_q(t)_{xz}|^2$, and
\begin{gather}
  \E_q\sup_t\bigl|\langle\psi_q(t),N_{\le m}\psi_q(t)\rangle-s_m\bigr|\le\SigF\le M_1(F),\qquad
  \sup_t\bigl|\langle\Psi(t),N_{\le m}\Psi(t)\rangle-s_m\bigr|\le\SigF,\label{eq:main-transport}\\
  \sup_t\bigl\langle\Psi(t),(N_{\le m}-s_m)^2\Psi(t)\bigr\rangle\le\ThF+\SigpF^2\le4M_2(F)+M'_1(F)^2.
  \label{eq:main-transport2}
\end{gather}
In particular $\Prob_{\Psi(t)}(|N_{\le m}-s_m|\ge k)\le(4M_2(F)+M'_1(F)^2)/k^2$ for all $t$ and $k>0$.
\item \emph{(Dressed correlations.)}  For Fock $\psi$ and $x\neq y$,
$\sup_t|\langle\Psi(t),\cdr_x^\dagger\cdr_y\Psi(t)\rangle|\le2\sum_{k\ge\lceil|x-y|/2\rceil}F(k)$; here
$\cdr_x^\dagger\cdr_y=f_x^\dagger(\prod_{x\le k<y}\sigma^z_{b_k})f_y$ for $x<y$ (\cref{thm:duality}(d)).
\end{enumerate}
All bounds are uniform in $N$, in the position of the cut and in $t$.
\end{theorem}

\begin{proof}
(a) By \eqref{eq:heisenberg}, $\cdr_x(t)=\sum_qP_q\sum_zU_q(t)_{xz}\cdr_z$.  The $\cdr$'s commute with every
$G_k$ (\cref{thm:duality}(d)), hence with $P_q$, and satisfy the CAR.  So
$\{\cdr_x(t),\cdr_y^\dagger\}=\sum_qP_q\sum_zU_q(t)_{xz}\delta_{zy}$, which is \eqref{eq:main-anticomm}.
Then $P_q\{\cdr_x(t),\cdr_y^\dagger\}P_q=U_q(t)_{xy}P_q$ with $P_q\neq0$, which gives the identity in
\eqref{eq:main-prop}; the bound is \eqref{eq:cons1}.  The last statement is \cref{lem:sectornorm}(c), together with
$\E_q|U_q(t)_{xy}|\le\E_q\sup_s|U_q(s)_{xy}|$.

(b) Here $\psi_q(t)=c^\dagger(U_q(t)\delta_y)\ket0$, so that
$\langle\psi_q(t),n_x\psi_q(t)\rangle=|U_q(t)_{xy}|^2$, and
$\langle\Psi(t),n_x\Psi(t)\rangle=\E_q|U_q(t)_{xy}|^2$ by \cref{cor:sector}(b),(c).  Exchange the finite sums,
use $\sup_t\E\le\E\sup_t$, and apply \eqref{eq:cons4}.

(c) By \cref{cor:sector}(c), $\langle\psi_q(t),n_x\psi_q(t)\rangle=\sum_{s\in S}|U_q(t)_{xs}|^2$.  Using
$\sum_{x=1}^N|U_q(t)_{xs}|^2=1$ for the $s\in S\cap[1,m]$,
\[
  \langle\psi_q(t),N_{\le m}\psi_q(t)\rangle-s_m=\sum_{x\le m,\ s\in S,\ s>m}|U_q(t)_{xs}|^2-\sum_{x>m,\ s\in S,\ s\le m}|U_q(t)_{xs}|^2 .
\]
Both sums are non-negative.  Since $|U_q(t)_{xs}|=|U_q(t)_{sx}|$, both are at most
\[
  D_q(t):=\sum_{x\le m<z}|U_q(t)_{xz}|^2,
\]
and the difference of two numbers in $[0,D]$ has modulus at most $D$.  Next, $D_q(t)\le Y_q:=\sum_{x\le m<z}Q_q(x,z)$, and $\E Y_q\le\SigF$ by \cref{lem:annealed}(ii).  The
LGT statement follows from \cref{cor:sector}(b), which gives
\[
  \langle\Psi(t),N_{\le m}\Psi(t)\rangle=\E_q\langle\psi_q(t),N_{\le m}\psi_q(t)\rangle .
\]
For \eqref{eq:main-transport2}: by \cref{cor:sector}(c) the law of $N_{\le m}$ in $\Psi(t)$ is the $q$-average of its
laws in the $\psi_q(t)$.  Hence
\[
  \langle\Psi(t),(N_{\le m}-s_m)^2\Psi(t)\rangle=\E_q\bigl[\Var_q(t)+(\langle\psi_q(t),N_{\le m}\psi_q(t)\rangle-s_m)^2\bigr],
\]
where $\Var_q(t)$ is the variance of $N_{\le m}$ in $\psi_q(t)$.  By \cref{lem:slater}(b) with $A=[1,m]$,
$\Var_q(t)=\sum_{x\le m<y}|K_q(t)_{xy}|^2\le\sum_{x\le m<y}|K_q(t)_{xy}|$, whose $\E\sup_t$ is at most $\ThF$ by
\cref{lem:annealed}(ii).  The second term is at most $\E Y_q^2\le\SigpF^2$.  The probability bound is
Chebyshev's inequality.

(d) $\langle\Psi(t),\cdr_x^\dagger\cdr_y\Psi(t)\rangle=\E_qK_q(t)_{yx}$ by \cref{cor:sector}(b),(c); apply
\cref{lem:annealed}(i).
\end{proof}

\begin{remark}[On the statement]\label{rem:main-transport}
(i) The supremum over $t$ is \emph{inside} the sector average, which is stronger than a bound on
$\sup_t\E_q$: in a typical sector the dressed fermion stays localized for all times.
(ii) By \cref{lem:productlaw}(c), \cref{thm:transport} holds verbatim for every $\sigma^z$-product gauge
state $\ket{s}_z$ in place of $\ket{\uparrow\cdots\uparrow}_z$.  For any gauge state whose sector law
has density at most $K$ with respect to the uniform law (e.g.\ $K=2$ for the open chain,
\cref{prop:variants}(i)), all bounds of \cref{thm:transport} hold with an extra factor $K$ (for Fock $\psi$ every
sector carries $\ket n$ up to a phase, by \cref{lem:productlaw}(b)).  This uses that they are upper bounds on
averages of non-negative sector functionals (\cref{rem:productlaws}(c)); lower bounds, such as
\cref{thm:imbalance}, do not transfer in this way.  For the dimer law of \cref{lem:productlaw}(e) and
$0<2|h|<J$, the moment bound (b) fails uniformly in $N$, and so does the analogue of
Hypothesis~\ref{hyp:EC} (\cref{thm:dimer}(d)).
(iii) Hypothesis~\ref{hyp:EC} enters only through \eqref{eq:hypEC} and its consequences
\eqref{eq:cons1}--\eqref{eq:cons4}:
(a) uses \eqref{eq:cons1}, (b) uses \eqref{eq:cons4}, (c) uses \eqref{eq:hypEC} itself together with
\eqref{eq:cons2} and \eqref{eq:cons3}, and (d) uses \eqref{eq:cons2}.  Part (c) is the gauge-theory analogue
of the disorder-averaged particle-transport bound for disordered XX chains in
\cite[Thm.~1.1]{ANSS2016}.
(iv) \emph{Numerical check} (\texttt{code/main\_checks.py}, part A; $N=60$, $\lambda=2$, CDW, $300$
sectors, $62$ times in $[0,10^5]$): the deterministic inequalities in the proof of (c) and the moment
inequalities hold to $4\times10^{-14}$, and $U_q(t)^\top=U_q(t)$ to $3\times10^{-16}$.  The measured
$\E\sup_t$ of the transported charge across the middle cut is $0.52$.  The bound $\SigF$ evaluated with
the empirical correlator $F_{\rm emp}$ (the non-increasing envelope of $\max_x\E Q_q(x,x+d)$, an
estimate, not a proof) is $20$.  For one particle, $\E\sup_t\sum_x(|x-y|+1)^2|U_q(t)_{xy}|^2=33$,
against $2M_2(F_{\rm emp})\approx5\times10^2$.
\end{remark}

\subsection{A sector-averaged zero-velocity Lieb--Robinson bound}\label{sec:main:LR}

For $X\subset[1,N]$ let $\Alg(X)$ be the $*$-algebra generated by $\{\cdr_x:x\in X\}$ and $\Alg^+(X)$ its even
part (spanned by monomials of even degree).  For $Y\subset[1,N]$ let $\BAlg^+(Y)$ be the $*$-algebra generated by
$\Alg^+(Y)$ and $G_1,\dots,G_N$.  All these operators are gauge invariant.  The generators of $\Alg^+(X)$ are local
in the original variables: for $x<y$, $\cdr_x^\dagger\cdr_y=f_x^\dagger(\prod_{x\le k<y}\sigma^z_{b_k})f_y$ and
$\cdr_x\cdr_y=f_xf_y\prod_{x\le k<y}\sigma^z_{b_k}$, and $\cdr_x^\dagger\cdr_x=n_x$ (\cref{thm:duality}(d)).
$\BAlg^+(Y)$ also contains the electric energies $\sigma^x_{b_{j-1}}\sigma^x_{b_j}=G_j(-1)^{n_j}$, $j\in Y$.  The
operators $\gamma_{2x-1}:=\cdr_x+\cdr_x^\dagger$ and $\gamma_{2x}:=i(\cdr_x^\dagger-\cdr_x)$, $x\in X$, are self-adjoint
unitaries.  For an ordered set $\mathcal S$ of indices let $\gamma_{\mathcal S}$ be the ordered product.  Every
$A\in\Alg(X)$ has a unique expansion $A=\sum_{\mathcal S}a_{\mathcal S}\gamma_{\mathcal S}$ (see the proof), and we
put $|A|_1:=\sum_{\mathcal S}|a_{\mathcal S}|\le4^{|X|}\norm{A}$.

\begin{proposition}[Sector-averaged zero-velocity Lieb--Robinson bound]\label{prop:LR}
Let $X,Y\subset[1,N]$ and $B\in\BAlg^+(Y)$.
\begin{enumerate}[label=(\alph*)]
\item \emph{(Sector-wise bound; no hypothesis.)}  For a monomial $A=\cdr^{\sharp}_{x_1}\cdots\cdr^{\sharp}_{x_k}$
($x_i\in X$, each $\cdr^\sharp$ equal to $\cdr$ or $\cdr^\dagger$), and for every $q$ and $t$,
\begin{equation}\label{eq:main-LRsector}
  \norm{P_q[A(t),B]}\le2\norm B\sum_{i=1}^k\sum_{y\in Y}|U_q(t)_{x_iy}| .
\end{equation}
For general $A\in\Alg(X)$, $\norm{P_q[A(t),B]}\le8|A|_1\norm B\sum_{x\in X,y\in Y}|U_q(t)_{xy}|$.
\item \emph{(Sector-averaged zero-velocity bound.)}  Under Hypothesis~\ref{hyp:EC},
\begin{equation}\label{eq:main-LR}
\begin{gathered}
  \E_q\sup_{t\in\R}\norm{P_q[A(t),B]P_q}\le2\norm B\sum_{i=1}^k\sum_{y\in Y}F(|x_i-y|),\\
  \text{resp.}\quad\E_q\sup_{t\in\R}\norm{P_q[A(t),B]P_q}\le8|A|_1\norm B\sum_{x\in X,y\in Y}F(|x-y|).
\end{gathered}
\end{equation}
Consequently, let $\rho$ be a density matrix whose sector weights satisfy $\Tr(\rho P_q)\le K2^{-N}$ for all
$q$.  Then $\sup_t|\Tr(\rho[A(t),B])|$ is at most $K$ times the right side of \eqref{eq:main-LR}.  This covers $K=1$ and
$\rho=\ket{\Psi_0}\bra{\Psi_0}$, $\Psi_0=\ket{\uparrow\cdots\uparrow}_z\otimes\psi$, with \emph{arbitrary} $\psi$, as well
as $\rho=4^{-N}\one$.
\item \emph{(The operator norm sees the clean sector.)}  For all $N$, $x,y\in[1,N]$ and $t$,
\begin{equation}\label{eq:main-LRnorm}
\begin{gathered}
  \norm{\{\cdr_x(t),\cdr_y^\dagger\}}=\norm{[\cdr_x(t),n_y]}=\max_q|U_q(t)_{xy}|,\\
  \norm{[n_x(t),n_y]}=\max_q|U_q(t)_{xy}|\bigl(1-|U_q(t)_{xy}|^2\bigr)^{1/2}.
\end{gathered}
\end{equation}
The maximum includes the clean sector $q=\Zone=(1,\dots,1)$.  It has probability $2^{-N}$ under the sector
law, which vanishes as $N\to\infty$.  In particular, for $N\ge2$,
\[
  \sup_t\norm{\{\cdr_1(t),\cdr_N^\dagger\}}\ge\Bigl(\frac{3}{2(N+1)}\Bigr)^{1/2},\qquad
  \sup_t\norm{[n_1(t),n_N]}\ge\Bigl(\frac{3}{4(N+1)}\Bigr)^{1/2},
\]
whereas the sector average in \eqref{eq:main-prop} is at most $F(N-1)$.
\item \emph{(No operator-norm zero-velocity bound.)}  There is no function $G:\N_0\to[0,\infty)$ such that
$\sum_dd^2G(d)^2<\infty$ and $\sup_t\norm{[n_x(t),n_y]}\le G(|x-y|)$ for all $N\ge1$ and all $x,y\in[1,N]$.
The same holds with $[n_x(t),n_y]$ replaced by $\{\cdr_x(t),\cdr_y^\dagger\}$ or by $[\cdr_x(t),n_y]$.
\end{enumerate}
\end{proposition}

\begin{proof}
(a) All operators involved are gauge invariant, so they commute with $P_q$.  Replacing each $G_j$ in $B$ by
$q_j$ gives $B_q\in\Alg^+(Y)$ with $P_qB=P_qB_q$.  If $z\notin Y$, then $\cdr_z$ anticommutes with $\cdr_y$ and
$\cdr_y^\dagger$ for all $y\in Y$, hence commutes with every even monomial, and $[\cdr_z,B_q]=0$.  By
\eqref{eq:heisenberg},
\[
  P_q[\cdr_x(t),B]=\sum_{z}U_q(t)_{xz}P_q[\cdr_z,B_q]=\sum_{y\in Y}U_q(t)_{xy}P_q[\cdr_y,B_q],
\]
and $\norm{P_q[\cdr_y,B_q]}\le\norm{\cdr_yP_qB_q}+\norm{P_qB_q\cdr_y}\le2\norm{P_qB}\le2\norm B$, because $P_q$
commutes with $\cdr_y$.  The same holds for $\cdr_x(t)^\dagger$, with $\overline{U_q(t)_{xy}}$.  For a monomial,
$[a_1\cdots a_k,B]=\sum_ia_1\cdots a_{i-1}[a_i,B]a_{i+1}\cdots a_k$.  Moving $P_q$ through
$a_1(t),\dots,a_{i-1}(t)$, which have norm one, gives \eqref{eq:main-LRsector}.

For general $A$ we first justify the expansion.  The products $\gamma_{\mathcal S}$, with $\mathcal S$ ranging over
the subsets of the $2|X|$ indices in increasing order, span $\Alg(X)$, since $\gamma_j^2=\one$ and distinct
$\gamma$'s anticommute.  Let $\tau:=4^{-N}\Tr$.  If $\mathcal S\neq\emptyset$, then $\tau(\gamma_{\mathcal S})=0$.
Indeed, for $|\mathcal S|$ odd, $\gamma_{\mathcal S}$ anticommutes with $(-1)^{N_f}$, since the $\cdr$'s are odd.  For
$|\mathcal S|$ even and $j\in\mathcal S$, $\gamma_j\gamma_{\mathcal S}\gamma_j=-\gamma_{\mathcal S}$.  In both cases
cyclicity of the trace gives $\tau(\gamma_{\mathcal S})=-\tau(\gamma_{\mathcal S})$.  Since
$\gamma_{\mathcal S}^*\gamma_{\mathcal S'}=\pm\gamma_{\mathcal S\triangle\mathcal S'}$, the $\gamma_{\mathcal S}$ are orthonormal for
$(A_1,A_2)\mapsto\tau(A_1^*A_2)$, the expansion is unique, and $a_{\mathcal S}=\tau(\gamma_{\mathcal S}^*A)$, so
$|a_{\mathcal S}|\le\norm A$.  As above,
$\norm{P_q[\gamma_{2x-1}(t),B]},\norm{P_q[\gamma_{2x}(t),B]}\le4\norm B\sum_{y\in Y}|U_q(t)_{xy}|$, and each
$x\in X$ carries two indices.  So
$\norm{P_q[\gamma_{\mathcal S}(t),B]}\le8\norm B\sum_{x\in X,y\in Y}|U_q(t)_{xy}|$.  Sum over $\mathcal S$.

(b) Take $\sup_t$ and $\E_q$ in (a) and use \eqref{eq:cons1}.  The statement on states is
\cref{lem:sectornorm}(b) applied to the gauge-invariant operator $X=[A(t),B]$.

(c) The first norm is $\max_q|U_q(t)_{xy}|$ by \eqref{eq:main-anticomm} and \cref{lem:sectornorm}(a).  Since
$[\cdr_z,\cdr_y^\dagger\cdr_y]=\delta_{zy}\cdr_y$, the computation in (a) gives
$P_q[\cdr_x(t),n_y]=U_q(t)_{xy}P_q\cdr_y$.  In the dual picture $P_q\cdr_y$ is $c_y\otimes\ket q\bra q$, of norm one.
For $[n_x(t),n_y]$ we use the dual picture in the sector $q$.  With $u:=U_q(t)$ and $w:=u^*\delta_x$, we have
$\Gamma(u)^*\dGamma(\ket{\delta_x}\bra{\delta_x})\Gamma(u)=\dGamma(\ket w\bra w)$, and
$[\dGamma(X_1),\dGamma(X_2)]=\dGamma([X_1,X_2])$.  Hence the operator is $\dGamma(Z)\otimes\ket q\bra q$ with
$Z=[\ket w\bra w,\ket{\delta_y}\bra{\delta_y}]=a\ket w\bra{\delta_y}-\bar a\ket{\delta_y}\bra w$, where
$a:=\langle w,\delta_y\rangle=u_{xy}$.  The matrix $iZ$ is Hermitian, of rank at most two and trace zero, so
its eigenvalues are $\pm\mu$ and $0$.  Here $2\mu^2=\Tr Z^*Z=2|a|^2(1-|a|^2)$, using $\norm w=1$.  For a
Hermitian one-body matrix with eigenvalues $(\lambda_i)$, the eigenvalues of its second quantization are the
sums $\sum_{i\in T}\lambda_i$; hence $\norm{\dGamma(Z)}=\mu$.

For the lower bound take $q=\Zone$.  The eigenvectors of $h_{\Zone}=-J\Delta_N+2h$ are
$\varphi_k(x)=\sqrt{2/(N+1)}\sin(\pi kx/(N+1))$, and $\varphi_k(N)^2=\varphi_k(1)^2$.  By \cref{lem:consequences}(c),
\[
  \sup_t|U_{\Zone}(t)_{1N}|^2\ge R_{\Zone}(1,N)=\sum_{k=1}^N\varphi_k(1)^4=\frac4{(N+1)^2}\sum_{k=0}^{N}\sin^4\frac{\pi k}{N+1}
  =\frac{3}{2(N+1)},
\]
because $\sin^4\theta=\frac38-\frac12\cos2\theta+\frac18\cos4\theta$, and
$\sum_{k=0}^{N}\cos\frac{2\pi jk}{N+1}=0$ for $j=1,2$ when $N+1\ge3$.  Put $s_0:=(3/(2(N+1)))^{1/2}\le2^{-1/2}$.
Since $t\mapsto|U_{\Zone}(t)_{1N}|$ is continuous, vanishes at $t=0$ and has supremum $\ge s_0$, it takes every value
$s<s_0$.  With $g(s):=s(1-s^2)^{1/2}\ge s/\sqrt2$ on $[0,2^{-1/2}]$, \eqref{eq:main-LRnorm} gives
$\sup_t\norm{[n_1(t),n_N]}\ge g(s_0)\ge s_0/\sqrt2$.

(d) Suppose $G$ exists.  Then $G(d)\to0$, so there is $d_0$ with $G(d)<\frac12$ for $d\ge d_0$.  Fix $N$ and
$x,y$ with $|x-y|\ge d_0$, and let $a(t):=U_{\Zone}(t)_{xy}$.  By \eqref{eq:main-LRnorm},
$g(|a(t)|)\le G(|x-y|)<\frac12$ for all $t$, where $g(s):=s(1-s^2)^{1/2}$.  The set
$T_0:=\{t:|a(t)|\le2^{-1/2}\}$ is closed and contains $t=0$.  On $[0,2^{-1/2}]$, $g(s)\ge s/\sqrt2$, so
$|a(t)|\le\sqrt2G(|x-y|)<2^{-1/2}$ on $T_0$, and by continuity $T_0$ is also open.  Hence $T_0=\R$ and
$|a(t)|\le\sqrt2G(|x-y|)$ for all $t$.  (For the other two operators, $|a(t)|\le G(|x-y|)$ directly by
\eqref{eq:main-LRnorm}.)  Now take $N=2L+1$ and $y=L+1$, and translate $[1,N]$ to $[-L,L]$.  By
\cref{lem:propagation}(b) with $v\equiv2h$, $\Lambda=[-L,L]$, $\Lambda'=\Z$, for fixed $t$ and $r$ the entry
$U_{\Zone}(t)_{L+1+r,L+1}$ converges as $L\to\infty$ to $(e^{-ith_\Z})_{r0}$, where $h_\Z:=-J\Delta_\Z+2h$.  So
$|(e^{-ith_\Z})_{r0}|\le\sqrt2G(|r|)$ for $|r|\ge d_0$ and all $t$, and
\[
  \sum_{r\in\Z}r^2|(e^{-ith_\Z})_{r0}|^2\le\sum_{|r|<d_0}r^2+4\sum_{d\ge d_0}d^2G(d)^2<\infty
\]
uniformly in $t$.  This contradicts the ballistic identity $\sum_rr^2|(e^{-ith_\Z})_{r0}|^2=2J^2t^2$ of
\cref{rem:gaussproj}(iii).
\end{proof}

\begin{remark}[Interpretation]\label{rem:main-LR}
\Cref{prop:LR}(b) is a \emph{state-dependent} zero-velocity bound.  It holds for every initial state whose
gauge field is $z$-polarized (or a $\sigma^z$-product state), whatever the matter state, and for the
infinite-temperature state, because these states give every sector the same weight.  Parts (c)--(d) show that
this restriction is necessary: the operator norm is a maximum over sectors, and the clean sector $q=\Zone$ (no
background charges, the gauge-invariant sector of the usual LGT formulation; \cref{rem:gaussproj}) is ballistic.
Even the on-site densities $n_1$, $n_N$, which are strictly local in the original variables, have a commutator
of norm $\ge(3/(4(N+1)))^{1/2}$ at suitable times, which is not exponentially small in the distance $N-1$.
Part (b) is the analogue, with the disorder average replaced by the sector average, of the disorder-averaged
zero-velocity Lieb--Robinson bound of Hamza, Sims and Stolz for the disordered XY chain \cite{HSS2012}.  The
difference is that here the randomness is not in the Hamiltonian but in the initial state, which is why the
bound is state dependent.
\emph{Numerical check} (\texttt{code/main\_checks.py}, parts B, C): for $N=4$ the identities
\eqref{eq:main-anticomm} and \eqref{eq:main-LRnorm} hold to $10^{-14}$ in exact diagonalization of the full
LGT.  The sector-wise bound \eqref{eq:main-LRsector} holds for $N=60$, $\lambda=2$, $200$ sectors and
$t\le10^5$; there $\E_q\sup_t\norm{P_q[A(t),n_{41}]}=4.0\times10^{-3}$ for
$A=\cdr_{11}^\dagger\cdr_{12}+{\rm h.c.}$.  Moreover $R_{\Zone}(1,N)=3/(2(N+1))$ to six digits for
$N=20,\dots,200$.
\end{remark}

\subsection{Infinite-time memory of a density wave}\label{sec:main:imbalance}

Let $N$ be even, $\psi$ the charge-density wave ($n_x=1$ iff $x$ is even), $\Gamma:=\diag((-1)^x)_{x=1}^N$,
$\Psi_0=\ket{\uparrow\cdots\uparrow}_z\otimes\psi$, and
\[
  I(t):=\frac2N\sum_{x=1}^N(-1)^x\langle\Psi(t),n_x\Psi(t)\rangle,\qquad I(0)=1 .
\]
For a sector $q$ let $(E_\alpha,\varphi_\alpha)$ be the eigenpairs of $\hq$ (real, orthonormal; simple spectrum),
$a_\alpha:=\langle\varphi_\alpha,\Gamma\varphi_\alpha\rangle$ and $V_q:=\diag(q)$.  Let
\begin{equation}\label{eq:main-kappa}
  \kappa_N:=\max_{B}\ \E_q\sum_{x\in B,\ y\in[1,N]\setminus B}R_q(x,y),
\end{equation}
the maximum over \emph{all} intervals $B\subset[1,N]$.

\begin{theorem}[Positive infinite-time imbalance for every $h\neq0$]\label{thm:imbalance}
Let $N$ be even.
\begin{enumerate}[label=(\roman*)]
\item For every $h\in\R$, the limit $\Ibar:=\lim_{T\to\infty}T^{-1}\int_0^TI(t)\,dt$ exists and
$\Ibar=\E_q\,N^{-1}\sum_\alpha a_\alpha^2\in[0,1]$.
\item \emph{(Anticommutator identity.)}  For every $h\in\R$, $q$ and $\alpha$:
$E_\alpha a_\alpha=2h\,\langle\varphi_\alpha,V_q\Gamma\varphi_\alpha\rangle$.  If $h=0$, then $a_\alpha=0$ for all
$\alpha$ and $\Ibar=0$ (chiral obstruction).
\item \emph{(Unconditional bounds.)}  Let $h\neq0$.  For every block length $1\le L\le N$,
\begin{equation}\label{eq:main-Ibound}
  \Ibar\ \ge\ \frac{\lfloor N/L\rfloor}{N}\cdot\frac{h^2\bigl(\sqrt{L/2}-2\kappa_N\bigr)_+^2}{L\,(J+|h|)^2},
  \qquad\text{and}\qquad \Ibar\ \ge\ \frac{h^2}{N(J+|h|)^2}>0 .
\end{equation}
\item \emph{(Uniform in $N$.)}  Under Hypothesis~\ref{hyp:EC}, $\kappa_N\le\kappa:=2\SigF\le2M_1(F)$ for all
$N$, and for all even $N\ge2$
\begin{equation}\label{eq:main-Iuniform}
  \Ibar\ \ge\ \frac{h^2}{16\,L_*\,(J+|h|)^2},\qquad L_*:=\max\{2,\lceil32\kappa^2\rceil\}.
\end{equation}
\end{enumerate}
\end{theorem}

\begin{proof}
(i) By \cref{cor:sector}(c), $\langle\psi_q(t),n_x\psi_q(t)\rangle=\sum_{s\in S}|U_{xs}|^2$ with $U=U_q(t)$ and
$S$ the even sites, and $P_S=\frac12(\one+\Gamma)$.  Hence
\begin{align*}
  I_q(t)&:=\frac2N\sum_x(-1)^x\langle\psi_q(t),n_x\psi_q(t)\rangle=\frac2N\Tr(P_SU^*\Gamma U)\\
  &=\frac1N\Tr\Gamma+\frac1N\Tr(\Gamma U^*\Gamma U)=\frac1N\Tr(\Gamma U^*\Gamma U),
\end{align*}
since $\Tr\Gamma=0$ for even $N$.  In the eigenbasis,
$\Tr(\Gamma U^*\Gamma U)=\sum_{\alpha,\beta}e^{it(E_\beta-E_\alpha)}|\langle\varphi_\alpha,\Gamma\varphi_\beta\rangle|^2$, and
$T^{-1}\int_0^Te^{i\omega t}dt\to\one[\omega=0]$.  The spectrum is simple (\cref{lem:consequences}(a)), so the
time average of $I_q$ is $N^{-1}\sum_\alpha a_\alpha^2$.  Since $I(t)=\E_qI_q(t)$ (\cref{cor:sector}(b)) is a finite
average, (i) follows; $|a_\alpha|\le1$.

(ii) $\Delta_N$ connects only sites of opposite parity, so $\Gamma\Delta_N\Gamma=-\Delta_N$, and
$\Gamma\hq+\hq\Gamma=2h(\Gamma V_q+V_q\Gamma)=4hV_q\Gamma$.  Taking the expectation in $\varphi_\alpha$ gives
$2E_\alpha a_\alpha=4h\langle\varphi_\alpha,V_q\Gamma\varphi_\alpha\rangle$.  If $h=0$, then $\hq=-J\Delta_N$, whose
eigenvalues $-2J\cos(\pi k/(N+1))$ do not vanish for even $N$; so $a_\alpha=0$.

(iii) Fix $q$ and an interval $B$ with $|B|=L$.  Put $m_\alpha(B):=\sum_{x\in B}\varphi_\alpha(x)^2$,
$r_x:=(-1)^xq_x$ and $R(x,B):=\sum_{y\in B}R_q(x,y)$.

\emph{Step 1.}  By (ii), $E_\alpha a_\alpha=2h\sum_xr_x\varphi_\alpha(x)^2$.  Multiplying by $m_\alpha(B)$ and
summing over $\alpha$ gives
\[
  X_B:=\sum_\alpha E_\alpha a_\alpha m_\alpha(B)=2h\sum_xr_xR(x,B).
\]
\emph{Step 2 (Cauchy--Schwarz).}
\begin{gather*}
  X_B^2\le\Bigl(\sum_\alpha E_\alpha^2m_\alpha(B)\Bigr)\Bigl(\sum_\alpha a_\alpha^2m_\alpha(B)\Bigr),\\
  \sum_\alpha E_\alpha^2m_\alpha(B)=\sum_{x\in B}\norm{\hq\delta_x}^2\le L(2J^2+4h^2)\le4L(J+|h|)^2 .
\end{gather*}

\emph{Step 3 (pointwise lower bound).}  Write
$X_B/(2h)=\sum_{x\in B}r_x+\sum_xr_x(R(x,B)-\one_B(x))$.  By \eqref{eq:Rdef}, $\sum_yR_q(x,y)=1$ and $R_q$ is
symmetric, so $\sum_x|R(x,B)-\one_B(x)|=\sum_{x\in B}R(x,B^c)+\sum_{x\notin B}R(x,B)=2\sum_{x\in B,y\notin B}R_q(x,y)$.
Hence
\[
  |X_B|\ge2|h|\Bigl(\Bigl|\sum_{x\in B}r_x\Bigr|-2\sum_{x\in B,y\notin B}R_q(x,y)\Bigr).
\]
\emph{Step 4 (averaging).}  The $(r_x)_{x\in B}$ are i.i.d.\ uniform signs.  By Szarek's sharp Khintchine
inequality \cite{Szarek1976}, $\E|\sum_{x\in B}r_x|\ge(L/2)^{1/2}$.  Hence
$\E|X_B|\ge2|h|(\sqrt{L/2}-2\kappa_N)$ and $\E X_B^2\ge(\E|X_B|)^2\ge4h^2(\sqrt{L/2}-2\kappa_N)_+^2$.

\emph{Step 5 (blocks).}  Let $B_j:=[(j-1)L+1,jL]$ for $1\le j\le\lfloor N/L\rfloor$.  Since
$\sum_jm_\alpha(B_j)\le\norm{\varphi_\alpha}^2=1$, Steps 2 and 4 give
\[
  \E\sum_\alpha a_\alpha^2\ge\sum_j\E\sum_\alpha a_\alpha^2m_\alpha(B_j)\ge\sum_j\frac{\E X_{B_j}^2}{4L(J+|h|)^2}
  \ge\Bigl\lfloor\frac NL\Bigr\rfloor\frac{h^2(\sqrt{L/2}-2\kappa_N)_+^2}{L(J+|h|)^2}.
\]
Divide by $N$ and use (i).

\emph{Step 6 (global bound).}  For $B=[1,N]$, $R(x,B)=1$, so $X_{[1,N]}=2h\sum_xr_x$ exactly and
$\E X_{[1,N]}^2=4h^2N$, while $\sum_\alpha E_\alpha^2=\Tr\hq^2\le4N(J+|h|)^2$.  Step 2 gives
$\E\sum_\alpha a_\alpha^2\ge h^2/(J+|h|)^2$.

(iv) $\E R_q(x,y)\le F(|x-y|)$ by \eqref{eq:cons1}, and an interval has at most $2d$ crossing pairs at distance
$d$; so $\kappa_N\le2\SigF$.  Let $L:=L_*$.  Then $4\kappa\le\sqrt{L/2}$, so
$(\sqrt{L/2}-2\kappa)^2\ge L/8$.  If $N\ge2L$, then $\lfloor N/L\rfloor\ge N/L-1\ge N/(2L)$, and
\eqref{eq:main-Ibound} gives $\Ibar\ge h^2/(16L(J+|h|)^2)$.  If $N<2L$, the global bound gives
$\Ibar\ge h^2/(N(J+|h|)^2)>h^2/(2L(J+|h|)^2)$.
\end{proof}

\begin{corollary}[Imbalance on the open chain]\label{cor:imbalanceopen}
Consider the open chain of \cref{prop:variants}(i), i.e.\ the model of \cite{SKKM2017} without the boundary link,
with the density-wave matter state and $z$-polarized links, and let $\Ibar^{\rm open}$ be its infinite-time
imbalance.  Under Hypothesis~\ref{hyp:EC}, for all even $N\ge4$,
\[
  \Ibar^{\rm open}\ \ge\ \frac{h^2}{16\,L'_*\,(J+|h|)^2},\qquad L'_*:=\max\{2,\lceil128\kappa^2\rceil\},
\]
with $\kappa$ as in \cref{thm:imbalance}(iv).
\end{corollary}

\begin{proof}
By \cref{prop:variants}(i) the sector law is uniform on the parity hyperplane $\{q:\prod_jq_j=\pi\}$,
$\pi=(-1)^{N/2}$, and the proof of \cref{thm:imbalance}(i) gives $\Ibar^{\rm open}=\E'_qN^{-1}\sum_\alpha a_\alpha^2$,
where $\E'$ is the average over this hyperplane.  Three properties of $\E'$ replace the independence of the
charges:
(a) for every set $B$ with $|B|\le N-1$ the charges $(q_x)_{x\in B}$ are i.i.d.\ uniform under $\E'$, since for
$x_0\notin B$ the restriction to $[1,N]\setminus\{x_0\}$ is a bijection from the hyperplane onto $\{\pm1\}^{N-1}$;
(b) $\E'f\le2\,\E f$ for every $f\ge0$;
(c) for $N\ge3$, $\E'(\sum_xr_x)^2=N$, by (a) applied to pairs of sites.
Steps~1--3 of the proof of \cref{thm:imbalance} are deterministic.  In Step~4, (a) gives Szarek's bound for
$|B|=L\le N-1$, and by (b) the quantity $\kappa'_N$ defined with $\E'$ satisfies $\kappa'_N\le2\kappa_N\le2\kappa$.
Step~5 is used with blocks of length $L\le N/2$, and Step~6 with (c).  The argument of
\cref{thm:imbalance}(iv), with $\kappa$ replaced by $2\kappa$, gives the claim.
\end{proof}

\begin{remark}[On \cref{thm:imbalance}]\label{rem:main-imbalance}
(i) \emph{Constants.}  The factor $(J+|h|)^2$ comes from $\norm{\hq}\le2J+2|h|$; the proof gives the slightly
better $\frac12(J^2+2h^2)$ in its place.  Writing $J+h$ instead of $J+|h|$ would be wrong for $h<0$: at
$h=-0.9J$, $N=40$, the ``global bound'' would be $2.03>1$.  The sharp constant $2^{-1/2}$ is Szarek's theorem.
An elementary alternative is H\"older's inequality,
$(\E S^2)^3\le(\E|S|)^2\,\E S^4$ for $S=\sum_{x\in B}r_x$, with $\E S^2=L$ and $\E S^4=3L^2-2L$.  It gives
$\E|S|\ge(L/3)^{1/2}$, and then \eqref{eq:main-Ibound} holds with $\sqrt{L/3}$ in place of $\sqrt{L/2}$, and
\eqref{eq:main-Iuniform} with $24$ and $48$ in place of $16$ and $32$.
(ii) \emph{Mechanism.}  The identity (ii) shows that a nonzero diagonal $a_\alpha$ requires $h\neq0$: for $h=0$ the
chiral symmetry $\Gamma\hq\Gamma=-\hq$ pairs the eigenvalues $\pm E$ and forces $a_\alpha=0$.  In the clean sector
$q=\Zone$, $(E_\alpha-2h)a_\alpha=0$ and again $\Ibar_{\Zone}=0$ (\cref{rem:gaussproj}(ii)).  Positivity of $\Ibar$ is
produced by the sector average, through the anticoncentration of $\sum_{x\in B}r_x$.
(iii) \emph{The time average needs no $\sup_t$ localization.}  The proof uses only the decay of $\E R_q$, a
time-averaged, eigenfunction-level quantity.  A lower bound on $\inf_tI(t)$ is not proved here.
(iv) \emph{Infinite volume.}  We do not formulate \cref{thm:imbalance} in the infinite-volume dynamics of
\cref{thm:infvol}: $\Ibar$ is an infinite-time average of a volume average, and an infinite-volume version
would require exchanging $T\to\infty$ and $N\to\infty$.  The natural thermodynamic-limit statement is the
uniformity in $N$ of \eqref{eq:main-Iuniform}.  It relies on Hypothesis~\ref{hyp:EC} on finite intervals,
i.e.\ on Theorem~\ref{thm:EC}.  We do not know how to obtain it from the Ge--Zhao bound on $\Z$ alone:
$\kappa_N$ involves the infinite-time quantities $R_q$ of the finite chain, which finite propagation speed
does not relate to their counterparts on $\Z$.
(v) \emph{Numerical check} (\texttt{code/main\_checks.py}, part D).  The identity (ii) and Steps 1--3 hold to
$10^{-13}$.  For $N=8$, $h=0.7J$ and all $256$ sectors, the exact time average up to $T=10^7$ agrees with
$\E_qN^{-1}\sum_\alpha a_\alpha^2=0.190469$ to $10^{-8}$.  With $\kappa_N$ computed as the maximum over all
intervals ($300$ sectors), $N=64$: $\lambda=0.5,1,2,4,-1.8$ give $\Ibar=0.023,0.082,0.173,0.247,0.158$,
$\kappa_N=9.0,4.6,2.4,2.1,2.5$, and the best of the two bounds \eqref{eq:main-Ibound} is $6\times10^{-4}$,
$1.7\times10^{-3}$, $3.9\times10^{-3}$, $6.9\times10^{-3}$ and $3.5\times10^{-3}$, respectively.  So the
rigorous bounds are valid but small.  For $\lambda\le0.5$ the convergence in $N$ is slow: at $\lambda=0.5$,
$\Ibar=0.0241$ for $N=50$ ($400$ sectors), and $\Ibar=0.01968(14)$, $0.01732(9)$, $0.01596(6)$ for $N=100,200,400$
($2000$ sectors; \cref{tab:numerics-imbalance}).  See also \cref{fig:imbalance}.
\end{remark}

\subsection{Entanglement}\label{sec:main:entanglement}

\begin{theorem}[Uniform area laws]\label{thm:arealaw}
Let $\psi=\ket n$ be a Fock state and $1\le\ell<N$.
\begin{enumerate}[label=(\alph*)]
\item \emph{(Deterministic chain.)}  For every $q$ and $t$, with $T_q(t):=P_{[\ell+1,N]}K_q(t)P_{[1,\ell]}$,
\[
  S(\rho^q_A(t))\le2\ln2\,\norm{T_q(t)}_1\le2\ln2\sum_{x\le\ell<y}|K_q(t)_{xy}| .
\]
\item \emph{(Sector-averaged area law.)}  Under Hypothesis~\ref{hyp:EC},
\[
  \E_q\sup_{t\in\R}S(\rho^q_A(t))\le2\ln2\,\ThF\le8\ln2\,M_2(F),\qquad\text{hence}\qquad
  \sup_t\SG(t)\le8\ln2\,M_2(F),
\]
uniformly in $N$, $\ell$ and the Fock state $\psi$.
\item \emph{(Definite-charge initial states.)}  Let $\Psi_0=\ket{\xi}_\sigma\otimes\ket n$ with
$\xi\in\{\pm1\}^N$ (links in $\sigma^x$-eigenstates), and $q=q(n,\xi)$ as in \cref{lem:productlaw}(d).  Then
$\Sfull(t)=S(\rho^q_A(t))$ for all $t$.  If $\xi$ is uniformly distributed, so is $q$, and under
Hypothesis~\ref{hyp:EC}
\[
  \E_\xi\sup_t\Sfull(t)\le8\ln2\,M_2(F),\qquad
  \Prob_\xi\Bigl(\sup_t\Sfull(t)\ge s\Bigr)\le\frac{8\ln2\,M_2(F)}{s}\quad(s>0),
\]
and $\E_\xi\Sfull(t)=\SG(t)$ for every $t$.
\end{enumerate}
\end{theorem}

\begin{proof}
(a) $\psi_q(t)=\prod_{s\in S}c^\dagger(U_q(t)\delta_s)\ket0$ is the Slater determinant of $W=U_q(t)\ell^2(S)$,
with projection $K_q(t)$.  By \cref{lem:slater}(a), $S(\rho^q_A(t))=\sum_i\bent(\nu_i)$.  By \cref{lem:binent}
this is at most $2\ln2\sum_i\sqrt{\nu_i(1-\nu_i)}$, and \cref{lem:slater}(b) finishes the proof.  (b)
\cref{lem:annealed}(ii) and \eqref{eq:main-constants}.  (c) The identity is \eqref{eq:quenchedSfull}, which rests
on \cref{lem:cutlocal}(c).  For fixed $n$, $\xi\mapsto q(n,\xi)$ is a bijection of $\{\pm1\}^N$
(\cref{lem:productlaw}(b)), so it maps the uniform law to the uniform law.  The bounds follow from (b) and
Markov's inequality.
\end{proof}

\begin{remark}\label{rem:main-arealaw}
(i) \emph{Method.}  The argument of (a)--(b) is that of Abdul-Rahman and Stolz \cite{AbdulRahmanStolz2015} and of
Abdul-Rahman, Nachtergaele, Sims and Stolz \cite[proof of Thm.~1.2]{ANSS2016} for disordered XY chains: the
binary-entropy bound of \cref{lem:binent} (inequality (61) of \cite{ANSS2016}), control of the entropy by the
entries of the correlation matrix across the cut, and an averaged eigenfunction-correlator bound.  Here the
average is over the sectors, not over a random Hamiltonian.
(ii) For an interval $A=[a,b]$ in the bulk the same argument, with \cref{lem:annealed}(iii), bounds the
sector-averaged entropy of the $c$-modes in $A$ by $4\ln2\,\ThF$ (two cuts).  The same counting, on $\Z$ and with
$\FZ$ in place of $F$, gives \cref{thm:infvol}(d).
(iii) \emph{Numerical check} (\texttt{code/main\_checks.py}, part E).  On a grid of $2\times10^6$ points
$\bent(\nu)\le2\ln2\sqrt{\nu(1-\nu)}$ holds; $d_1=0.153$ and $d_k<0$ for $2\le k\le2\times10^5$.
\cref{lem:slater} is confirmed to $10^{-15}$ against many-body Slater determinants ($8$ modes; left, interior and
non-contiguous $A$).  The chain in (a) is confirmed along the dynamics ($N=60$, $\lambda=2$, $t\le10^5$).
For definite-charge states the identity in (c) is confirmed by exact diagonalization of the full LGT ($N=4,5$)
to $4\times10^{-14}$.  The measured $\E\sup_tS(\rho^q_A(t))$ for the half chain ($N=60$, $\lambda=2$, CDW) is
$1.7$.  The bound $2\ln2\,\Theta(F_{\rm emp})\approx1.3\times10^3$ is valid but crude.
\end{remark}

We now compare the entropies of the superposed state.  Let $\psi=\ket n$ be a Fock state,
$\Psi_0=\ket{\uparrow\cdots\uparrow}_z\otimes\psi$, $\Phi(t):=U^\dagger\Psi(t)=2^{-N/2}\sum_qe^{iht\sum_jq_j}\psi_q(t)\otimes\ket q$
(\cref{cor:sector}(a)), and $N_A:=\sum_{x\le\ell}n_x$.  Define
\begin{itemize}
\item $\Sfull(t):=S_A(\Psi(t))$, the matter-plus-gauge entanglement of $A$ (\S\ref{sec:model:cut});
\item $\SG(t):=\E_qS(\rho^q_A(t))$, the entanglement conditional on all Gauss charges;
\item $\Smix(t):=S(\rho_{\rm mix}(t))$, $\rho_{\rm mix}(t):=\E_q\rho^q_A(t)$;
\item $\Sfm(t):=S(\rho^f_A(t))$, where $\rho^f_A(t)$ is the reduced state of $\Psi(t)$ on the bare fermion modes
$f_1,\dots,f_\ell$, with \emph{all} links traced out;
\item $\SN(t):=\Hsh(p_t)$ with $p_t(k):=\Prob_{\Psi(t)}(N_A=k)$; and $p^q_t(k)$ the law of $N_A$ in $\psi_q(t)$;
\item $H_t:=\Hsh(\mu_t)$, where $\mu_t$ is the law of the random density matrix $q\mapsto\rho^q_A(t)$ (a law on a finite
set).
\end{itemize}

\begin{proposition}[Entropy taxonomy]\label{prop:entropies}
For every $t\in\R$:
\begin{enumerate}[label=(\alph*)]
\item \emph{(Duality defect.)}  $|\Sfull(t)-S_{A'}(\Phi(t))|\le\ln2$.
\item \emph{(Conditional versus mixed.)}  $\rho_{\rm mix}(t)$ is the marginal of $\rho_{A'}(\Phi(t))$ on the modes
$c_1,\dots,c_\ell$, and
\[
  \SG(t)\le\Smix(t)\le\min\{\SG(t)+H_t,\ \ell\ln2\},\qquad H_t\le N\ln2,\qquad H_0=0 .
\]
\item \emph{(Number entropy as a lower bound.)}  $p_t=\E_qp^q_t$, and for $X\in\{\mathrm{full},\mathrm{mix},f\}$
\[
  S_X(t)=\SN(t)+\sum_kp_t(k)\,S\bigl(\rho_X^{(k)}(t)\bigr)\ \ge\ \SN(t),
\]
where $p_t(k)\rho^{(k)}_X(t)$ is the block of the respective reduced state on $\{N_A=k\}$.  For the conditional
entropy one only has $\SG(t)\ge\E_q\Hsh(p^q_t)$, and $\E_q\Hsh(p^q_t)\le\SN(t)\le\E_q\Hsh(p^q_t)+H_t$.
\item \emph{(Superposition versus ensemble.)}  Let $\rho^{\rm ens}:=\sum_qP_q\ket{\Psi_0}\bra{\Psi_0}P_q$.  Then
\[
  \rho^{\rm ens}=2^{-N}\sum_{\xi\in\{\pm1\}^N}\ket{\xi}\bra{\xi}_\sigma\otimes\ket n\bra n=2^{-N}\one_{\rm links}\otimes\ket n\bra n,
\]
and for all gauge-invariant $O_1,\dots,O_m$ and $t_1,\dots,t_m\in\R$,
\begin{equation}\label{eq:main-ens}
  \langle\Psi_0,O_1(t_1)\cdots O_m(t_m)\Psi_0\rangle=\Tr\bigl(\rho^{\rm ens}O_1(t_1)\cdots O_m(t_m)\bigr).
\end{equation}
In particular, the following are the same for $\Psi(t)$ and for the ensemble
$\rho^{\rm ens}(t):=e^{-itH}\rho^{\rm ens}e^{itH}$: the law of every gauge-invariant observable; the restriction of
the state to the dressed-fermion algebra $\Alg([1,\ell])$, which is represented through $U$ by $\rho_{\rm mix}(t)$;
hence $\Smix(t)$; and $\SN(t)$.  The members $\ket\xi_\sigma\otimes\ket n$ of the
ensemble are definite-charge states, and their average full entanglement is $\SG(t)$
(\cref{thm:arealaw}(c)).
\end{enumerate}
\end{proposition}

\begin{proof}
(a) \cref{lem:cutlocal}(b) with $\Phi=\Phi(t)$.

(b) Tracing $\ket{\Phi(t)}\bra{\Phi(t)}$ over the $c$-modes $\ell+1,\dots,N$ and over all $\tau$'s kills the
off-diagonal terms in $q$, because the $\ket q$ are orthonormal.  What remains is
$2^{-N}\sum_q\Tr_{c_{>\ell}}\ket{\psi_q(t)}\bra{\psi_q(t)}=\rho_{\rm mix}(t)$.  The first inequality is concavity of $S$.
Group the sectors according to the value of $\rho^q_A(t)$: $\rho_{\rm mix}(t)=\sum_\sigma\mu_t(\sigma)\sigma$.
The argument in the proof of \cref{lem:pinching} gives $S(\sum_i\mu_i\sigma_i)\le\sum_i\mu_iS(\sigma_i)+\Hsh(\mu)$:
$\mu_i\sigma_i\le\sum_j\mu_j\sigma_j$ and operator monotonicity of $\ln$ give
$-\Tr\sigma_i\ln\sum_j\mu_j\sigma_j\le S(\sigma_i)-\ln\mu_i$, which we average over $i$.  This gives
$\Smix\le\SG+H_t$.  Also $\Smix\le\ln\dim\cF_{\le\ell}=\ell\ln2$, and $H_t\le\Hsh(q)=N\ln2$.  At $t=0$,
$\rho^q_A(0)$ does not depend on $q$, so $H_0=0$.

(c) $p_t=\E_qp^q_t$ is \cref{cor:sector}(c).  Let $\cH=\cH_1\otimes\cH_2$ with $N_f=N_1\otimes\one+\one\otimes N_2$, both
with integer spectrum, and let $N_f\Psi=M\Psi$.  Then $\Psi=\sum_k(Q_k\otimes Q'_{M-k})\Psi$, with $Q_k,Q'_j$ the
spectral projections.  The partial trace over $\cH_2$ of the cross terms $k\neq k'$ vanishes, because
$Q'_{M-k}Q'_{M-k'}=0$.  Hence $\Tr_2\ket\Psi\bra\Psi=\bigoplus_kp(k)\rho^{(k)}$ with
$p(k)=\langle\Psi,(Q_k\otimes\one)\Psi\rangle$, and the entropy of such a block-diagonal state is
$\Hsh(p)+\sum_kp(k)S(\rho^{(k)})$.  Apply this to $\Psi(t)$, which has $N_f\Psi(t)=|S|\Psi(t)$:
\begin{itemize}
\item for $\Sfull$, with $\cH_1=\cH_A$, where the links carry no particle number;
\item for $\Sfm$, with $\cH_1=\cF_{\le\ell}$ and $\cH_2=\cF_{>\ell}\otimes(\C^2)^{\otimes N}$;
\item for each $\psi_q(t)$, with $\cH_1=\cF^c_{\le\ell}$.  This shows that every $\rho^q_A(t)$ is block diagonal
with weights $p^q_t$, hence so is $\rho_{\rm mix}(t)$, with weights $\E_qp^q_t=p_t$.
\end{itemize}
In all cases the weights are those of $N_A$, by \cref{cor:sector}(c).  The same per-sector decomposition gives
$S(\rho^q_A(t))\ge\Hsh(p^q_t)$.  The inequality $\E_q\Hsh(p^q_t)\le\Hsh(\E_qp^q_t)$ is concavity, and the upper
bound is the classical case of the mixing inequality in (b), after grouping the sectors by the value of
$\rho^q_A(t)$, which determines $p^q_t$.

(d) By \cref{cor:sector}(a) and \eqref{eq:U},
$P_q\Psi_0=2^{-N/2}U(\ket n\otimes\ket q)=2^{-N/2}\ket n\otimes\ket{x(n,q)}_\sigma$.  As $q$ runs over $\{\pm1\}^N$,
$x(n,q)$ runs over $\{\pm1\}^N$ (\cref{lem:gauss}(iii)), which gives the formula for $\rho^{\rm ens}$.  The operator
$Y:=O_1(t_1)\cdots O_m(t_m)$ is gauge invariant, because $H$ commutes with every $G_j$.  So
\[
  \langle\Psi_0,Y\Psi_0\rangle=\sum_{q,q'}\langle P_{q'}\Psi_0,YP_q\Psi_0\rangle=\sum_q\langle P_q\Psi_0,YP_q\Psi_0\rangle
  =\Tr(\rho^{\rm ens}Y).
\]
For the next statement: every element of the CAR algebra of $c_1,\dots,c_\ell$ corresponds under $U$ to an element
of $\Alg([1,\ell])$, which is gauge invariant.  So $\rho_{\rm mix}(t)$ is determined by gauge-invariant expectation
values; $\SN$ is the entropy of the law of the gauge-invariant $N_A$.  The last statement is
\eqref{eq:quenchedSfull} for the members of the ensemble.
\end{proof}

\begin{remark}[$\SG$ can be smaller than $\SN$]\label{rem:main-SGSN}
The inequality $\SG\ge\SN$ is \emph{false} in general.  The number entropy is the entropy of the
\emph{mixture} $p_t=\E_qp^q_t$, while $\SG$ only dominates the \emph{average} $\E_q\Hsh(p^q_t)$, and concavity
of $\Hsh$ goes the wrong way.  Physically, different charge configurations shift the mean of $N_A$ by
different amounts.  This spread contributes to $\SN$, but not to the entanglement in any single sector.  Exact
numerics (\texttt{code/main\_checks.py}, part F; $N=10$, $J=h=1$, $\ell=5$, CDW, all $1024$ sectors) give
$\SG<\SN$ at early times:
\begin{center}
\begin{tabular}{c|cccccc}
$t$ & $\Sfull$ & $S_{A'}(\Phi)$ & $\SG$ & $\Smix$ & $\Sfm$ & $\SN$\\\hline
$0.5$ & $0.927$ & $0.476$ & $0.459$ & $1.391$ & $1.352$ & $0.461$\\
$1$ & $1.446$ & $0.975$ & $0.605$ & $2.652$ & $2.571$ & $0.670$\\
$2$ & $2.265$ & $1.760$ & $0.713$ & $2.709$ & $2.664$ & $0.800$\\
$10$ & $3.403$ & $2.919$ & $0.913$ & $2.566$ & $2.674$ & $0.816$\\
$10^2$ & $5.000$ & $4.786$ & $0.980$ & $2.775$ & $2.826$ & $0.890$\\
$10^5$ & $5.382$ & $5.331$ & $0.945$ & $2.902$ & $2.929$ & $0.876$
\end{tabular}
\end{center}
(The maxima are $2\ell\ln2=6.93$ for $\Sfull$ and $\ell\ln2=3.47$ for $\Smix$, $\Sfm$.)  All other inequalities of
\cref{prop:entropies} hold in these data, with $H_t\approx6.93\approx N\ln2$ for $t>0$ (the $1024$ sectors give
$1022$ or $1023$ distinct $\rho^q_A(t)$).  The pipeline is cross-checked against exact diagonalization of the full LGT for $N=4$, to
$10^{-15}$.  See also \cref{fig:entropies}.
\end{remark}

\begin{remark}[Superposition versus sampling; relation to experiment]\label{rem:superposition-main}
\Cref{prop:entropies}(d) makes exact, for the present model, the statement that no gauge-invariant measurement,
at any number of times, distinguishes the coherent superposition over charge sectors from the classical
ensemble of definite-charge states.  The ensemble is the infinite-temperature state of the gauge field.  Only
quantities that are not expectation values of gauge-invariant observables distinguish the two:
\begin{itemize}
\item gauge-variant observables, such as the Loschmidt-type $\langle\sigma^z_{b_j}(t)\rangle$ of
\eqref{eq:loschmidt}, which vanishes identically in the ensemble;
\item entanglement functionals.  The superposed state has $\Sfull$, which numerically keeps growing (table
above, Conjecture~\ref{conj:logent}), while the ensemble members have, on average, the bounded $\SG$ (\cref{thm:arealaw}).
\end{itemize}
Gyawali et al.\ \cite{Gyawali2024} observed the same dichotomy on a quantum processor, for a $\Z_2$ LGT in one
and two spatial dimensions prepared in superpositions over gauge sectors: local observables show localization,
while R\'enyi entropy measurements distinguish superposition-prepared states from direct disorder sampling.
Their model is a Trotterized $\Z_2$ LGT with qubit matter; it is non-integrable and does not reduce to free
fermions, so it is not the model studied here.
\Cref{prop:entropies}(d) is an exact version of this dichotomy for the 1D model of \cite{SKKM2017}, stated for von
Neumann entropies.  We claim no more than this qualitative correspondence with the experiment.
\end{remark}

\subsection{Number entropy}\label{sec:main:number}

\begin{lemma}[Entropy of integer variables {\cite{Massey1988}}]\label{lem:intentropy}
Let $X$ be an integer-valued random variable taking finitely many values.  Then
$\Hsh(X)\le\frac12\ln\bigl(2\pi e(\Var X+\frac1{12})\bigr)$.
\end{lemma}

\begin{proof}
Let $U$ be uniform on $[0,1)$ and independent of $X$, and $Y:=X+U$.  $Y$ has the density
$f(y)=p(\lfloor y\rfloor)$, where $p$ is the law of $X$.  So its differential entropy is
$-\int f\ln f=-\sum_kp(k)\ln p(k)=\Hsh(X)$, and $\Var Y=\Var X+\frac1{12}=:\sigma^2$.  Let $\varphi$ be the
Gaussian density with the mean and variance of $Y$.  By Jensen's inequality,
$\int_{f>0}f\ln(\varphi/f)\le\ln\int_{f>0}\varphi\le0$, so
\[
  \Hsh(X)=-\int f\ln f\le-\int f\ln\varphi=\tfrac12\ln(2\pi\sigma^2)+\frac{\E(Y-\E Y)^2}{2\sigma^2}=\tfrac12\ln(2\pi e\sigma^2).
\]
This is the maximum-entropy property of the Gaussian; see \cite[Ch.~8]{CoverThomas2006}.
\end{proof}

\begin{theorem}[Bounded number entropy]\label{thm:numberentropy}
Assume Hypothesis~\ref{hyp:EC}.  Let $\psi$ be a Fock state, $\Psi_0=\ket{\uparrow\cdots\uparrow}_z\otimes\psi$ and
$A=[1,\ell]$, $1\le\ell<N$.  Then, for all $t\in\R$ and all $N$,
\begin{equation}\label{eq:main-SN}
  \SN(t)\le\tfrac12\ln\bigl(2\pi e\,(K_1+K_2+\tfrac1{12})\bigr),\qquad
  K_1:=\ThF\le4M_2(F),\quad K_2:=\SigpF^2\le M'_1(F)^2 .
\end{equation}
Moreover $\E_q\sup_t\Hsh(p^q_t)\le\frac12\ln\bigl(2\pi e(K_1+\frac1{12})\bigr)$ for the number entropy in a single sector.
\end{theorem}

\begin{proof}
Let $s_\ell:=|S\cap A|$.  Since $p_t$ is the law of $N_A=N_{\le\ell}$ in $\Psi(t)$,
$\Var_{p_t}(N_A)\le\langle\Psi(t),(N_A-s_\ell)^2\Psi(t)\rangle\le K_1+K_2$ by \eqref{eq:main-transport2}.  Apply
\cref{lem:intentropy}.  For the second statement, \cref{lem:intentropy} in each sector gives
$\sup_t\Hsh(p^q_t)\le\frac12\ln(2\pi e(\sup_t\Var_q(t)+\frac1{12}))$.  By Jensen's inequality for the concave
logarithm and the proof of \cref{thm:transport}(c), the $q$-average is at most
$\frac12\ln(2\pi e(\E\sup_t\Var_q(t)+\frac1{12}))\le\frac12\ln(2\pi e(K_1+\frac1{12}))$.
\end{proof}

\begin{remark}[Size of the constant]\label{rem:main-numberentropy}
\emph{Numerical check} (\texttt{code/main\_checks.py}, part G; $N=64$, $\lambda=2$, CDW, $A=[1,32]$, $400$
sectors, $t\le10^5$).  The measured number entropy stays below $1.03$ for all sampled times, and
$\Var_{p_t}(N_A)\le0.46$.  The intermediate quantities are $\E\sup_t\Var_q=0.54$ ($\le K_1$) and
$\E\sup_t(\langle N_A\rangle_q-s_\ell)^2=0.34$ ($\le K_2$).  The theorem's bound \eqref{eq:main-SN}, evaluated with
the empirical correlator, is about $6.0$.  With the exponential majorant $F(d)=\min\{1,e^{-0.18d}\}$ of the
data it is $6.3$.  These are illustrations: \eqref{eq:main-SN} holds with the true $F$ of
Hypothesis~\ref{hyp:EC}, whose constants are not explicit.  The value $1.11$ obtained by inserting the
\emph{measured} variance into \cref{lem:intentropy} is not a consequence of the theorem.  Lemma~\ref{lem:intentropy}
itself was tested on $3000$ random integer laws and on binomial laws.
\end{remark}

\subsection{Thermodynamic limit}\label{sec:main:infvol}

The chain $[1,N]$ has a fixed left end, and the region $A=[1,\ell]$ touches it (\cref{rem:boundary}).  Bounds on
$\ell^2(\Z)$ therefore do not directly control the finite-volume quantities above.  In this subsection we use
two-sided chains, observables in the bulk, and \emph{only} the Ge--Zhao bound \eqref{eq:GeZhao}, in the form of
\cref{prop:Zline}.

\paragraph{Two-sided chains.}
For $M\in\N$ let $\Lambda_M:=[-M,M]\cap\Z$, with sites $j\in\Lambda_M$, links $b_j=(j,j+1)$ for $j\in\Lambda_M$, the
boundary link $b_M$, and $\sigma^x_{b_{-M-1}}:=\one$.  Let $H^{(M)}$, $G_j$ and $\cdr^{(M)}_j=f_j\prod_{k=j}^M\sigma^z_{b_k}$ be
defined as in \cref{def:model} and \eqref{eq:dressed}.  The relabelling $j\mapsto j+M+1$ identifies this with the model
on $[1,2M+1]$, so \cref{sec:model} and the preceding results apply.  Fix $n\in\{0,1\}^\Z$ with occupied set $S$, let
$\Psi^{(M)}_0:=\ket{\uparrow\cdots\uparrow}_z\otimes\ket{n|_{\Lambda_M}}$, and $\Psi^{(M)}(t):=e^{-itH^{(M)}}\Psi^{(M)}_0$.  We
realize all sector laws on one probability space: $q=(q_x)_{x\in\Z}$ is i.i.d.\ uniform on $\{\pm1\}$, $\E$ is the
corresponding expectation, and the sector of $\Lambda_M$ is $q|_{\Lambda_M}$.  Let
\begin{gather*}
  h^{(M)}_q:=-J\Delta_{\Lambda_M}+2h\diag(q|_{\Lambda_M}),\qquad h^\Z_q:=-J\Delta_\Z+2h\diag(q)\ \text{on }\ell^2(\Z),\\
  U^{(M)}_q(t):=e^{-ith^{(M)}_q},\qquad U^\Z_q(t):=e^{-ith^\Z_q},
\end{gather*}
and $K^\Z_q(t):=U^\Z_q(t)P_SU^\Z_q(t)^*$, a projection on $\ell^2(\Z)$.  Let $Q^\Z_q(x,y)$ be the correlator of
\cref{prop:Zline} for $A_\Z=h^\Z_q$.  By \cref{lem:normalization} and \eqref{eq:GeZhao},
$\E\,Q^\Z_q(x,y)\le\FZ(|x-y|)$ with
\[
  \FZ(d):=\min\{1,C_\Z e^{-\eta_\Z d}\},
\]
where $\eta_\Z\in(0,\gamma^\E)$ and $C_\Z$ are as in \cref{prop:Zline}.  $\FZ$ is non-increasing, and all its moments
$M_p(\FZ)$, $M'_p(\FZ)$ are finite.  (Lemma~\ref{lem:normalization} is stated on finite intervals; on $\Z$ the same
diagonal unitary conjugates $h^\Z_q$ to $J(\Delta_\Z+\lambda\diag(q))$.)  By \cref{lem:propagation}(b), the entries of
$U^\Z_q(t)$ are uniform limits of the entries of $U^{(M)}_q(t)$, which depend on finitely many $q_x$; hence
$U^\Z_q(t)_{xy}$ and $K^\Z_q(t)_{xy}$ (an absolutely convergent series of products of such entries) are measurable
in $q$.  By the
proof of \cref{cor:ECinfvol}, $Q^\Z_q(x,y)$ is the supremum of $|f(h^\Z_q)_{xy}|$ over a countable dense subset of the
continuous $f$ with $|f|\le1$ on $[-2J-2|h|,2J+2|h|]$, and each $f(h^\Z_q)_{xy}$ is a uniform limit of polynomial
entries $p(h^\Z_q)_{xy}$; so $Q^\Z_q(x,y)$ is measurable too.  All functions of $t$ below are continuous, so the suprema
over $t$ may be taken over rational $t$.

\paragraph{Local gauge-invariant observables.}
For a finite interval $X=[a,b]\subset\Z$ and $M$ with $X\subset[-M+1,M-1]$, let $\GAlg(X)$ be the $*$-algebra generated
by $\cdr^{(M)\sharp}_x\cdr^{(M)\sharp}_y$ ($x,y\in X$) and $G_j$ ($j\in X$).  By \cref{thm:duality}(d) its generators are
$f^\sharp_xf^\sharp_y\prod_{x\le k<y}\sigma^z_{b_k}$, $n_x$ and $\sigma^x_{b_{j-1}}(-1)^{n_j}\sigma^x_{b_j}$.  These act only on
the sites of $X$ and the links $b_{a-1},\dots,b_b$, and, as even elements, they do not depend on $M$.  So
$\GAlg(X)$ is a fixed algebra of local gauge-invariant observables.  It contains the dressed hoppings, the densities,
the charges and the electric energies in $X$.  For $O\in\GAlg(X)$, $U^\dagger OU=\sum_qO_q\otimes\ket q\bra q$, where
$O_q$, in the even CAR algebra of the $c$-modes in $X$, is obtained by $\cdr\mapsto c$, $G_j\mapsto q_j$; it depends
on $q$ only through $q|_X$, and $\norm{O_q}\le\norm O$.  We regard $O_q$ as an operator on $\cF_X$.

\begin{theorem}[Thermodynamic limit; unconditional]\label{thm:infvol}
Let $\lambda\neq0$.  In the setting above:
\begin{enumerate}[label=(\alph*)]
\item \emph{(Limit dynamics.)}  For every finite interval $X$, $O\in\GAlg(X)$ and $t\in\R$, the limit
\[
  \omega_t(O):=\lim_{M\to\infty}\langle\Psi^{(M)}(t),O\,\Psi^{(M)}(t)\rangle=\E\,\Tr_{\cF_X}\bigl(\rhoqf{K^\Z_q(t)|_X}\,O_q\bigr)
\]
exists, where $K|_X:=P_XKP_X$ on $\ell^2(X)$.  The convergence holds sector by sector, uniformly in $q$, and
uniformly for $t$ in compact sets.
\item \emph{(Propagator, moments, transport.)}  For all $x,y\in\Z$, $\E\sup_t|U^\Z_q(t)_{xy}|\le\FZ(|x-y|)$, and
$\E|U^{(M)}_q(t)_{xy}|\to\E|U^\Z_q(t)_{xy}|$ for every $t$.  If $S=\{y\}$ (one particle), then
$\langle\Psi^{(M)}(t),n_x\Psi^{(M)}(t)\rangle\to\E|U^\Z_q(t)_{xy}|^2$, and for $0\le p\le2$
\[
  \sup_t\sum_{x\in\Z}(|x-y|+1)^p\,\E|U^\Z_q(t)_{xy}|^2\le2M_p(\FZ).
\]  For general $S$ and every $m\in\Z$, the
transported charge
\[
  \mathcal T_m(t):=\lim_{M\to\infty}\Bigl(\langle\Psi^{(M)}(t),N_{[-M,m]}\Psi^{(M)}(t)\rangle-|S\cap[-M,m]|\Bigr)
\]
exists for every $t$, and $\sup_t|\mathcal T_m(t)|\le\Sigma_1(\FZ)$.
\item \emph{(Lieb--Robinson.)}  Let $X\subset\Z$ be finite and $Y\subset\Z$ a finite interval.  Let $\Alg^+(X)$ be the
even part of the algebra generated by the $\cdr^{(M)}_x$, $x\in X$; it is spanned by products of the even monomials
$\cdr^{(M)\sharp}_x\cdr^{(M)\sharp}_y$, which do not depend on $M$.  Then for $A\in\Alg^+(X)$, $B\in\GAlg(Y)$ and every $t$,
\[
  \limsup_{M\to\infty}\E\norm{P_q[A(t),B]}\le8|A|_1\norm B\sum_{x\in X,y\in Y}\FZ(|x-y|),
\]
and likewise with the monomial bound of \eqref{eq:main-LR}.  The right side does not depend on $t$.
\item \emph{(Area law for bulk intervals.)}  For a finite interval $A\subset\Z$, let $S^{(M)}_q(t)$ be the entropy of
the restriction of $\psi^{(M)}_q(t)$ to the $c$-modes in $A$ (\cref{lem:slater}), $\SG^{(M)}(t):=\E S^{(M)}_q(t)$ the
bulk conditional entanglement, and $S^\Z_q(t):=S(\rhoqf{K^\Z_q(t)|_A})$.  For every $t$, $S^{(M)}_q(t)\to S^\Z_q(t)$
uniformly in $q$, so $\SG^{(M)}(t)\to\SG^\Z(t):=\E S^\Z_q(t)$.  Moreover
\[
  \E\sup_tS^\Z_q(t)\le4\ln2\,\Theta(\FZ)\le16\ln2\,M_2(\FZ).
\]
For the half-infinite regions $A_M$ (sites $-M,\dots,m$ and links $b_{-M},\dots,b_m$) one has, for every $t$,
$\limsup_M\E S(\rho^{q}_{A_M}(t))\le2\ln2\,\Theta(\FZ)$.  The same bound holds for $\limsup_M\E_\xi\Sfull^{(M)}(t)$,
for definite-charge initial states $\ket{\xi}_\sigma\otimes\ket{n|_{\Lambda_M}}$ with $\xi$ uniform.
\item \emph{(Number entropy for bulk intervals.)}  For a finite interval $A\subset\Z$ and every $t$, the law of $N_A$
in $\Psi^{(M)}(t)$ converges to $p_t:=\E\,\mathrm{PB}(\nu(K^\Z_q(t)|_A))$, where $\mathrm{PB}(\nu)$ denotes the
Poisson-binomial law with the eigenvalues $\nu$ of $K^\Z_q(t)|_A$.  Hence $\SN^{(M)}(t)\to\SN(t):=\Hsh(p_t)$, and
\[
  \sup_t\SN(t)\le\tfrac12\ln\bigl(2\pi e\,(2\Theta(\FZ)+4\Sigma_1'(\FZ)^2+\tfrac1{12})\bigr).
\]
\end{enumerate}
In (b)--(e) the bounds concern the limit objects, for every $t$; we do not claim that the limits $M\to\infty$ and
$\sup_t$ can be interchanged.  (By Theorem~\ref{thm:EC}, which is not used here, the finite-volume bounds also hold
uniformly in $M$.)
\end{theorem}

\begin{proof}
Throughout, $c:=|t|(2J+2|h|)$ and $w(d):=e^cc^d/d!$, so $\sum_dd^kw(d)<\infty$ for all $k$.  By
\cref{lem:propagation}(a), $|U^{(M)}_q(t)_{xy}|,|U^\Z_q(t)_{xy}|\le w(|x-y|)$ for all $M$ and $q$.  By
\cref{lem:propagation}(b) with $\Lambda=\Lambda_M$, $\Lambda'=\Z$, for $x,y\in\Lambda_M$,
\begin{equation}\label{eq:main-conv}
  |U^{(M)}_q(t)_{xy}-U^\Z_q(t)_{xy}|\le2e^c\frac{c^{n_0}}{n_0!},\qquad n_0=(M+1-|x|)+(M+1-|y|),
\end{equation}
uniformly in $q$.

\emph{Step 1: convergence of correlation matrices.}  Fix $x,y$, and let
$K^{(M)}_q(t):=U^{(M)}_q(t)P_{S\cap\Lambda_M}U^{(M)}_q(t)^*$.  By Cauchy--Schwarz,
\[
  |K^{(M)}_q(t)_{xy}-K^\Z_q(t)_{xy}|
  \le\Bigl(\sum_{s\in\Lambda_M}|U^{(M)}_{xs}-U^\Z_{xs}|^2\Bigr)^{1/2}+\Bigl(\sum_{s\in\Lambda_M}|U^{(M)}_{ys}-U^\Z_{ys}|^2\Bigr)^{1/2}
  +\sum_{s\notin\Lambda_M}w(|x-s|)w(|y-s|),
\]
where we used $\sum_s|U_{xs}|^2=1$ for both unitaries.  By \eqref{eq:main-conv}, every term is bounded independently of
$q$ and tends to $0$: for $|x|,|y|\le M/2$ and $M\ge2c$, each summand in the first two sums is at most
$(2e^cc^{\lceil M/2\rceil}/\lceil M/2\rceil!)^2$, and there are $2M+1$ of them.  The bound is also uniform for $|t|$ in
compact sets.

\emph{(a).}  By \cref{cor:sector}(b) for $\Lambda_M$,
$\langle\Psi^{(M)}(t),O\Psi^{(M)}(t)\rangle=\E\langle\psi^{(M)}_q(t),O_q\psi^{(M)}_q(t)\rangle$.  The vector $\psi^{(M)}_q(t)$ is the
Slater determinant with projection $K^{(M)}_q(t)$, and $O_q$ lies in the CAR algebra of the modes in $X$.  So by
\cref{lem:slater}(a) the inner expectation is $\Tr(\rhoqf{K^{(M)}_q(t)|_X}O_q)$.  By Step 1 and the uniform continuity
of $C\mapsto\rhoqf C$ (\cref{lem:quasifree}(a)),
$\norm{\rhoqf{K^{(M)}_q(t)|_X}-\rhoqf{K^\Z_q(t)|_X}}_1\to0$ uniformly in $q$, and $\norm{O_q}\le\norm O$.

\emph{(b).}  $\sup_t|U^\Z_q(t)_{xy}|\le Q^\Z_q(x,y)$ (take $f(E)=e^{-itE}$), so
$\E\sup_t|U^\Z_q(t)_{xy}|\le\FZ(|x-y|)$.  The convergence of $\E|U^{(M)}_q(t)_{xy}|$ is \eqref{eq:main-conv}.  For one
particle, $\langle\Psi^{(M)}(t),n_x\Psi^{(M)}(t)\rangle=\E|U^{(M)}_q(t)_{xy}|^2$ (\cref{thm:transport}(b)), which converges.
For every $t$, by Tonelli,
$\sum_x(|x-y|+1)^p\E|U^\Z_q(t)_{xy}|^2\le\E\sum_x(|x-y|+1)^pQ^\Z_q(x,y)\le2M_p(\FZ)$, since $|U|^2\le|U|\le Q^\Z_q$ and
there are at most two $x$ at each distance from $y$.  For transport, the proof of \cref{thm:transport}(c) gives
\[
  \langle\Psi^{(M)}(t),N_{[-M,m]}\Psi^{(M)}(t)\rangle-|S\cap[-M,m]|=\E\bigl[a^{(M)}_q(t)-b^{(M)}_q(t)\bigr],
\]
with the sums over $x,s\in\Lambda_M$
\[
  a^{(M)}_q:=\sum_{x\le m<s,\ s\in S}|U^{(M)}_{q,xs}|^2,\qquad b^{(M)}_q:=\sum_{s\le m<x,\ s\in S}|U^{(M)}_{q,xs}|^2 .
\]  Each term converges uniformly in $q$ and is dominated by $w(|x-s|)^2$, and
$\sum_{x\le m<s}w(s-x)^2=\sum_dd\,w(d)^2<\infty$.  By dominated convergence for series, uniformly in $q$,
\[
  a^{(M)}_q\to a^\Z_q:=\sum_{x\le m<s,\ s\in S}|U^\Z_{q,xs}|^2,
\]
and similarly for $b$.  So
$\mathcal T_m(t)=\E(a^\Z_q-b^\Z_q)$.  As $h^\Z_q$ is real symmetric, $U^\Z_q(t)$ is symmetric, so
$|a^\Z_q-b^\Z_q|\le\sum_{x\le m<s}|U^\Z_{q,xs}|^2\le\sum_{x\le m<s}Q^\Z_q(x,s)$.  The expectation of the last sum
is at most $\sum_dd\,\FZ(d)=\Sigma_1(\FZ)$.

\emph{(c).}  For each $M$ large, (a) of \cref{prop:LR}, applied to the chain $\Lambda_M$ and with $\GAlg(Y)\subset\BAlg^+(Y)$,
gives $\E\norm{P_q[A(t),B]}\le8|A|_1\norm B\sum_{x\in X,y\in Y}\E|U^{(M)}_q(t)_{xy}|$; the expansion coefficients of the
even element $A$ do not depend on $M$.  Let $M\to\infty$ using (b).

\emph{(d).}  By \cref{lem:slater}(a), the entropy of the $c$-modes of $A$ in $\psi^{(M)}_q(t)$ is
$\sum_i\bent(\nu_i(K^{(M)}_q(t)|_A))$.  Eigenvalues of Hermitian matrices are $1$-Lipschitz in operator norm (Weyl), and
$\bent$ is uniformly continuous on $[0,1]$; Step 1 gives the convergence, uniformly in $q$.  For the bound, $C:=K^\Z_q(t)|_A$
satisfies $C-C^2=T^*T$ with $T:=P_{\Z\setminus A}K^\Z_q(t)P_A$, because $K^\Z_q(t)$ is a projection on $\ell^2(\Z)$.  Hence,
by \cref{lem:quasifree}(b) and \cref{lem:binent}, $S^\Z_q(t)\le2\ln2\norm T_1\le2\ln2\sum_{x\in A,y\notin A}|K^\Z_q(t)_{xy}|$.
Here $T$ has finite rank, and the last step is the rank-one decomposition, which converges in trace norm when the sum
is finite.  Next, $|K^\Z_q(t)_{xy}|\le\sum_{z\in\Z}Q^\Z_q(x,z)Q^\Z_q(y,z)$ for all $t$, and
$\E[Q^\Z_q(x,z)Q^\Z_q(y,z)]\le\FZ(\max\{|x-z|,|y-z|\})$ because $Q^\Z_q\le1$.  The counting of
\cref{lem:annealed}(i),(iii) gives $\E\sup_tS^\Z_q(t)\le2\ln2\cdot2\Theta(\FZ)$.

For $A_M$: by \cref{thm:arealaw}(a),
\[
  S(\rho^q_{A_M}(t))\le2\ln2\sum_{x\le m<y;\,x,y\in\Lambda_M}|K^{(M)}_q(t)_{xy}| .
\]
We have
$|K^{(M)}_{q,xy}|\le\sum_zw(|x-z|)w(|y-z|)\le e^{2c}\sum_{n\ge|x-y|}(2c)^n/n!=:k(|x-y|)$, since distinct $z$ give distinct
pairs $(|x-z|,|y-z|)$ with sum at least $|x-y|$.  Moreover $\sum_dd\,k(d)<\infty$.  By Step 1 and dominated convergence,
the right side converges, uniformly in $q$, to $2\ln2\sum_{x\le m<y}|K^\Z_q(t)_{xy}|$, whose expectation is at most
$2\ln2\,\Theta(\FZ)$ by the counting of \cref{lem:annealed}(ii).  For the definite-charge states use
\eqref{eq:quenchedSfull} in $\Lambda_M$ (where $A_M$ is a left segment) and the uniformity of $q|_{\Lambda_M}$ when $\xi$ is
uniform.

\emph{(e).}  By \cref{lem:slater}(a) and \cref{lem:quasifree}(b), the generating function of $N_A$ in $\psi^{(M)}_q(t)$ is
$\det(\one+(z-1)K^{(M)}_q(t)|_A)$.  Its coefficients are polynomials in the entries, so they converge uniformly in $q$
by Step 1.  By \cref{cor:sector}(c), $p^{(M)}_t=\E\,\mathrm{PB}(\nu(K^{(M)}_q(t)|_A))\to p_t$; $N_A$ takes the finitely
many values $0,\dots,|A|$, so $\SN^{(M)}(t)\to\Hsh(p_t)$.  For the bound, let $s_A:=|S\cap A|$.  The law of total
variance gives $\Var_{p_t}(N_A)\le\E[\Var^\Z_q+(\mu^\Z_q-s_A)^2]$, where $\Var^\Z_q=\Tr(C-C^2)$ and
$\mu^\Z_q=\Tr C=\sum_{x\in A}\sum_{s\in S}|U^\Z_{q,xs}|^2$, with $C=K^\Z_q(t)|_A$.  As in (d),
$\Var^\Z_q=\norm T_2^2\le\sum_{x\in A,y\notin A}|K^\Z_{q,xy}|$, with $\E\sup_t$ at most $2\Theta(\FZ)$.  Using
$\sum_{x\in\Z}|U^\Z_{q,xs}|^2=1$ (all series converge absolutely),
$\mu^\Z_q-s_A=\sum_{x\in A,\,s\in S\setminus A}|U^\Z_{q,xs}|^2-\sum_{x\notin A,\,s\in S\cap A}|U^\Z_{q,xs}|^2$.  As in
\cref{thm:transport}(c), its modulus is at most $Y_q:=\sum_{x\in A,s\notin A}Q^\Z_q(x,s)$.  By Cauchy--Schwarz and
$\E(Q^\Z_q)^2\le\E Q^\Z_q$,
$\E Y_q^2\le\bigl(\sum_{x\in A,s\notin A}\FZ(|x-s|)^{1/2}\bigr)^2\le4\Sigma'_1(\FZ)^2$.  Conclude with
\cref{lem:intentropy}.
\end{proof}

\begin{remark}[Scope of \cref{thm:infvol}]\label{rem:main-infvol}
(i) The dressed fermion $\cdr^{(M)}_x$ itself is not a local operator in the limit: its Wilson line runs to the boundary link
$b_M$.  This is why (c) is formulated for even $A$.  The anticommutator \eqref{eq:main-anticomm} is nevertheless
controlled through (b).
(ii) The half-infinite statement in (d) is the infinite-volume counterpart of \cref{thm:arealaw}(b),(c): the region
contains the left end of $\Lambda_M$ and has a single cut, at $m$.
(iii) The imbalance (\cref{thm:imbalance}) is not covered; see \cref{rem:main-imbalance}(iv).
(iv) \emph{Numerical check} (\texttt{code/main\_checks.py}, part H).  The bounds of \cref{lem:propagation} hold to
$10^{-14}$ for $t\le5$ and $M\le40$.  With coupled sectors ($2000$ samples), $\E|U^{(M)}_q(3)_{0,3}|=0.13685$ for $M=4$ and $0.14341$ for
$M=8,16,32,64$.
\end{remark}

\subsection{Beyond the area law: sensitivity to distant charges and conjectures}\label{sec:main:conj}

The superposed state $\Psi(t)$ is a coherent superposition of $2^N$ sector evolutions.  Its full entanglement $\Sfull$
is not controlled by \cref{thm:arealaw}: by \cref{prop:entropies}(a) it is, up to $\ln2$, the entanglement of
$\Phi(t)=2^{-N/2}\sum_qe^{iht\sum q}\psi_q(t)\otimes\ket q$ across $A'|B'$.  This entanglement is generated by the dependence
of the evolution near the cut on the distant charges.  The charges $\tau^x_j$ are exactly conserved local
integrals of motion.  The rigorous input is the following sensitivity estimate.

\begin{lemma}[Duhamel sensitivity to distant charges]\label{lem:duhamel}
Let $q,q'\in\{\pm1\}^N$ differ exactly on $D\subset[1,N]$, and let $z\in[1,N]$, $s\in\R$.  Then
\[
  \norm{(U_q(s)-U_{q'}(s))\delta_z}\le4|h|\int_0^{|s|}\norm{\one_DU_{q'}(\pm r)\delta_z}\,dr\le4|h|\,|s|\sum_{x\in D}Q_{q'}(x,z),
\]
with the sign of $s$ in $\pm r$.  Under Hypothesis~\ref{hyp:EC}, with $q^D$ the configuration $q$ flipped on $D$,
\[
  \E_q\sup_{|s|\le t}\norm{(U_q(s)-U_{q^D}(s))\delta_z}\le4|h|\,t\sum_{x\in D}F(|x-z|).
\]
In particular, if $\dist(z,D)\ge R$ and $F(d)=\min\{1,Ce^{-\eta d}\}$, the right side is at most
$8|h|tCe^{-\eta R}/(1-e^{-\eta})$.
\end{lemma}

\begin{proof}
Let $s\ge0$ (the case $s<0$ is identical).  $\frac{d}{dr}\bigl(U_q(-r)U_{q'}(r)\bigr)=iU_q(-r)(\hq-h_{q'})U_{q'}(r)$, so
$U_{q'}(s)-U_q(s)=iU_q(s)\int_0^sU_q(-r)(\hq-h_{q'})U_{q'}(r)\,dr$.  Here $\hq-h_{q'}=2h\diag(q-q')=4h\diag(q\one_D)$, so
$\norm{(\hq-h_{q'})v}=4|h|\norm{\one_Dv}$.  Since $U_q$ is unitary, the first inequality follows.  Then
$\norm{\one_DU_{q'}(r)\delta_z}\le\sum_{x\in D}|U_{q'}(r)_{xz}|\le\sum_{x\in D}Q_{q'}(x,z)$.  The map $q\mapsto q^D$ preserves the
uniform law, so $\E_qQ_{q^D}(x,z)=\E_qQ_q(x,z)\le F(|x-z|)$.  The last claim sums the geometric series over the at most two
sites at each distance $d\ge R$.
\end{proof}

So two sectors that differ only at distance $\ge R$ from a site evolve indistinguishably near that site up to times of
order $e^{\eta R}/|h|$.  Heuristically, at time $t$ only the charges within distance
$R(t)\simeq\eta^{-1}\ln(|h|t)$ of the cut have dephased the sector components of $\Phi(t)$ near the cut.  This
suggests logarithmic growth of $\Sfull$, as for l-bit models of many-body localization; the difference is that here
the l-bits (the charges) are exact and their effect is mediated by free fermions.

\begin{conjecture}[Logarithmic growth of the full entanglement]\label{conj:logent}
Let $\psi$ be the CDW, $\Psi_0=\ket{\uparrow\cdots\uparrow}_z\otimes\psi$ and $\ell=\lfloor N/2\rfloor$.  There are constants
$C_1,C_2<\infty$, depending only on $J$ and $h$, such that $\Sfull(t)\le C_1\ln(2+|t|)+C_2$ for all $N$ and $t$.  Moreover
$\lim_{t\to\infty}\liminf_{N\to\infty}\Sfull(t)=\infty$.
\end{conjecture}

\emph{What is missing.}  \Cref{lem:duhamel} controls one-body propagators.  Turning it into an entanglement bound uniform in
$N$ requires a decoupling argument for a coherent superposition of $2^N$ sectors.  Entropy continuity bounds
(Fannes--Audenaert) cost $\varepsilon\ln\dim$, where $\dim$ is exponential in $N$.  A telescoping argument over the
distance of the flipped charges, with summable errors, would be needed; we have not carried this out.  The area-law
arguments for disordered XY chains \cite{AbdulRahmanStolz2015,ANSS2016}, which we used for $\SG$, concern quasi-free
states and do not apply to this superposition.  The lower bound would require a quantitative non-degeneracy of the sector phases
$e^{iht\sum q}$ and of the overlaps $\langle\psi_{q'}(t),\psi_q(t)\rangle$.

\begin{conjecture}[Volume law for the mixed entropies]\label{conj:volmix}
There is $c=c(J,h)>0$ such that, for all $N\ge2\ell$ and the CDW initial state,
$\liminf_{T\to\infty}T^{-1}\int_0^T\Smix(t)\,dt\ge c\,\ell$, and likewise for $\Sfm$.
\end{conjecture}

Since $\Smix(0)=\SG(0)=0$ for Fock states (\cref{prop:entropies}(b)), a time average or a restriction to $t\ge t_0$ is
necessary.  By \cref{prop:entropies}(d), $\Smix$ is the same for the superposition and for the ensemble.  So
Conjecture~\ref{conj:volmix} concerns the classical randomness of the charges: $\Smix-\SG$ is the Holevo information of the ensemble
$\{2^{-N},\rho^q_A(t)\}$, i.e.\ the information that the dressed matter in $A$ carries about the charges.  It does not
concern entanglement of $\Psi(t)$.  At strong coupling, each resonant pair of sites contributes an $O(1)$ mixture between a
frozen and an oscillating configuration, independently across pairs, which suggests a volume law.

\begin{remark}[Numerics for the conjectures]\label{rem:main-conjnum}
In the exact data of \cref{rem:main-SGSN} ($N=10$, $J=h=1$, $\ell=5$), $\Sfull$ grows from $1.45$ at $t=1$ to $3.40$, $5.00$ and $5.38$
at $t=10$, $10^2$, $10^5$, against the maximum $6.93$.  This is consistent with logarithmic growth followed by
finite-size saturation, but $N=10$ cannot separate the two.  $\Smix/(\ell\ln2)\approx0.84$ at late times.  \Cref{lem:duhamel}
is confirmed numerically in \texttt{code/main\_checks.py}, part I.
\end{remark}

\begin{remark}[Summary of the numerical checks]\label{rem:main-checks}
All inequalities and identities of this section were tested by \texttt{code/main\_checks.py} (output
\texttt{code/main\_checks.out}).  It uses the independent exact-diagonalization builders of \texttt{code/model\_duality\_check.py}
for the full LGT ($N\le5$) and the one-body operators $\hq$ for larger $N$.  No violation was found beyond round-off
($\le10^{-12}$).  The numbers quoted in the remarks of this section come from that output.  They illustrate the theorems
but are not used in any proof.
\end{remark}
